% Автоматически перенесено из ИСТОЧНИК_ВЕРСИЯ_3.docx

\section*{Источниковедение}
\addcontentsline{toc}{section}{Источниковедение}

В настоящем разделе последовательно рассматриваются различные группы источников диссертационного исследования. Внимание сосредоточено на ключевых памятниках обычного права Каталонии: «Барселонских обычаях» и связанных с ними городских сводах --- «Обычаях Льейды», «Обычаях Тортосы», «Кутюмах Перпиньяна», «Обычаях Орты», «Обычаях Миравета» и др. --- которые были составлены в XII--XIV~вв. При анализе каждого из этих текстов рассматриваются вопросы происхождения, датировки, языка, рукописной традиции, истории издания и места памятника в правовой культуре Каталонии.

Затем предметом рассмотрения становится актовый материал, прежде всего, судебные акты и соглашения, опубликованные в современных изданиях, посвящённых урегулированию конфликта в Средневековой Каталонии с IX по XII~вв. (акты IX -- XI~вв. используются как вспомогательные). Здесь важны не только содержание самих документов, но и принципы формирования соответствующих сборников, степень их репрезентативности, а также те исследовательские возможности и ограничения, которые связаны с использованием уже отобранного и изданного материала. Наконец, в завершение раздела кратко характеризуются вспомогательные источники, необходимые для сравнительного анализа и для выявления более широкого интеллектуального и правового контекста. Такая структура позволяет представить источниковую базу не как простую совокупность разнородных текстов, а как внутренне связанную систему, в которой нормативные памятники, документы практики и сравнительный материал последовательно дополняют друг друга.

\subsection*{§1. «Барселонские обычаи»}
\addcontentsline{toc}{subsection}{§1. «Барселонские обычаи»}

«Барселонские обычаи» --- один из наиболее известных памятников каталонского права, обладающий сложной структурой и множеством редакционных слоёв. Он представляет собой собрание постановлений графского совета Барселоны, посвящённых различным аспектам частного, публичного, уголовного и процессуального права.

Из текста преамбулы «Барселонских обычаев» мы узнаём, что придворные каталонские легисты рассматривали данный компендиум как сборник правовых обычаев, которые существовали на территории Каталонии и были записаны при Рамоне Беренгере~I Старом\footnote{Usatges de Barcelona. El codi a mitjan segle XII / ed. J. Bastardas i Parerа. Barcelona: Fundació Noguera, 1984 (1991). Us.1 (1--2)--3(4). Далее ссылки даются исключительно на это издание в сокращённой форме: Usatges de Barcelona. Us. N (n).}. Однако тот факт, что эти нормы (как правило, неписаные) были подвергнуты записи графской властью, говорит ещё и о том, что данный текст имел публичный характер, являлся нормативно-правовым и декларировал стремление политической власти упорядочить принципы регулирования отношений в обществе.

«Барселонские Обычаи» были составлены на латыни, которая характеризуется высокой степенью архаизации, частым применением терминов и конструкций, заимствованных из римского права и выраженным влиянием старокаталанского. В XIII~в. «Барселонские обычаи» были переведены на старокаталанский язык.

В вопросе о месте происхождения источника зарубежные и отечественные исследователи высказывают единодушие и подчёркивают, что местом его составления была Барселона\footnote{\emph{ Balari y Jovany D.J.} Origenes Históricos de Cataluña. Barcelona: Hijos de J. Jepús, 1899. P.~454.}. С центрами, где производился перевод текста этого памятника на старокаталанский язык, такой определённости нет, однако, можно с уверенностью сказать, что одним из таких центров тогда также являлась Барселона.

Определение точного времени создания «Барселонских обычаев» практически всегда вызывало споры в научном сообществе\footnote{Ранняя история полемики о датировке составления «Барселонских обычаев» подробно отражена в специализированном разделе у П.~Боннасси: \emph{Bonnassie P. }La Catalogne du milieu du Xe siècle à la fin du XIe siècle. Croissance et mutations d’une société. Toulouse: Association des publications de l’Université de Toulouse-Le Mirai, 1976. T.~2. P.~711--714.}. Если исследователи XIX--начала XX~вв. почти единодушно, за исключением Дж.~Балари-и-Джовани, руководствовавшегося лингвистическими особенностями текста\footnote{\emph{Idem}\emph{. }Op. cit. P.~453--466.}, определяли дату первичной записи этого правового памятника 1060--1070~гг.\footnote{Такой же датировки придерживался ряд исследователей середины XX~в.:\emph{ Люблинская А.Д. }Источниковедение истории Средних веков. Л.: Изд-во Ленингр. ун-та, 1955. С.~233--234; \emph{Lewis A.R. }The development of southern French and Catalan society, 718--1050. Austin: University of Texas, 1965. P.~380.}, то со временем практически общепринятая датировка была подвергнута критике. Составители «Барселонских обычаев» уже во второй статье указывали на то, что этот свод обычного права был записан при Рамоне Беренгере I и его супруге Альмодис де ла Марш\footnote{Usatges de Barcelona. Us. 2 (1--2), 3 (4). P.~50--54.}. В~хронике «Деяниях графов Барселонских» упоминание о записи «Барселонских обычаев» было отнесено к 1068~г. и связано с визитом кардинала Гуго Кандида в Барселону\footnote{\textbf{«}\textbf{Деяния}\textbf{ }\textbf{графов}\textbf{ }\textbf{Барселонских}\textbf{» }(лат. Gesta comitum Barcinonensium) --- латинская династическая хроника, созданная в Рипольского монастыря в XII--XIV~вв. и посвящённая происхождению, преемственности и политическим деяниям графов Барселонских. Цит. по: \emph{Bonnassie P. }La Catalogne du milieu du Xe siècle à la fin du XIe siècle. Croissance et mutations d’une société. Toulouse: , 1976. T.~2. P.~711.}. Однако при чтении текста «Барселонских обычаев» обращает на себя внимание ряд противоречий:

\begin{itemize}
\item во-первых, титулатура Рамона Беренгера I, в которой появляется характеристика \emph{«старый}\emph{»} (лат. vetus)\footnote{Usatges de Barcelona. Us. 1 (1--2). P.~48--49.}. Вероятно, она могла быть использована уже в правление его потомков-тёзок для облегчения определения правителя;
\item во-вторых, отдельные лица, перечисленные в составе графской курии, якобы утвердивших текст «Барселонских обычаев»\footnote{Ibid. Us. 3 (4.1). P.~52--54.}, не могли присутствовать там одновременно. Причиной этому было длившееся более 30 лет правление Рамона Беренгера I. За это время состав придворной курии несколько раз менялся. Более того, некоторые из упомянутых в тексте персоналий вошли в состав совета лишь при Рамоне Беренгере III\footnote{\emph{Kagay~D.J}. Introduction // The Usatges of Barcelona: The Fundamental Law of Catalonia. Philadelphia: University of Pennsylvania Press, 1994. P.~34.};
\item в-третьих, единичные и зачастую сомнительные упоминания о «Барселонских обычаях» в грамотах, относящихся к XI~в.\footnote{Так в 552 документах IX--XI~вв., опубликованных в 37-м томе серии «Каталонские юридические тексты» (2018) встречается лишь одно упоминание о «Барселонских обычаях» без прямой цитаты: «Et hoc fuit iudicatum per melius directum per \textbf{usum}\textbf{ de }\textbf{Barchinona} et de Urgello sine enganno» (Hoc est iudicium quod fuit iudicatum ad Sanctam Ceciliam…~: [1071] // Justícia i resolució de conflictes a la Catalunya medieval. Col·lecció diplomàtica. Segles IX--XI / dir. \emph{J.M.~Salrach} i Marès, T.de Montagut i Estragués. Barcelona: Parlament de Catalunya; Departament de Justícia de la Generalitat de Catalunya, 2018. Doc. 392. P.~677). См. также: \emph{Mayer Olivé~M.} Algunas notas a propósito de la compilación de los “Usatges de Barcelona” // Estudios de Latín Medieval Hispánico: Actas del V Congreso Hispánico de Latín Medieval. Barcelona. 7--10 de septiembre de 2009. Barcelona, 2012. P.~1001.} вкупе с активным использованием прямых цитирований из «Вестготской правды»\footnote{В тех же 552 документах IX--XI~вв., опубликованных в 37-м томе серии «Каталонские юридические тексты» (2018) встречается 100 упоминаний, ссылок и цитирований текста с указанием конкретных статей и титулов «Вестготской правды»: Justícia i resolució de conflictes a la Catalunya medieval. Col·lecció diplomàtica. Segles IX--XI / dir. \emph{J.M.~Salrach} i Marès, T.de Montagut i Estragués. Barcelona: Parlament de Catalunya; Departament de Justícia de la Generalitat de Catalunya, 2018. P.~1057--1067. См. также: \emph{Zimmermann M.} L’usage du Droit wisigothique en Catalogne du IXe au XIIe siècle: approches d’une signification culturelle. // Mélanges de la Casa de Velázquez. Madrid, 1973. T.~9. P.~237.}.
\end{itemize}
Вышеуказанные обстоятельства, как и поздняя рукописная традиция, не позволяют отнести появление первой редакции текста «Барселонских обычаев» ко времени правления Рамона Беренгера I. Стоит оговориться, что мы имеем в виду именно время записи существовавших правовых установлений, подчёркивая, что правовые традиции, основанные на переосмыслении наследия «Вестготской правды» и принадлежавшие области неписаного права, могли существовать на территории Каталонии и ранее XII~в.

Эта оговорка дала возможность ряду исследователей выдвинуть гипотезу так называемого «первоначального ядра» (фр. \emph{“}\emph{noyau}\emph{ }\emph{primitif}\emph{”)}, состоящего из 56 статей и относящегося ко второй половине XI~в., по мнению автора этой гипотезы, Р.~д’Абадаля-и-Виньялиса\footnote{Цит. по: \emph{Bastardas i Parera J.} Introducció // Usatges de Barcelona. El codi a mitjan segle XII / ed. J. Bastardas i Parerа. Barcelona: Fundació Noguera, 1984 (1991). P.~11.}. Частично такая трактовка композиции «Барселонских обычаев» основывалась на исследовании порядка записи обычаев в рукописной традиции, проведённым Г.М.~Брока-и-де-Монтагутом\footnote{\emph{Brocà i de Montagut G.M. }de. Els Usatges de Barcelona // Anuario del Instituto de Estudios Catalanes. Barcelona, 1914. P.~357--389.}.

Существование «первоначального ядра» отрицал П.~Боннасси, который, проанализировав текст этих 56 статей, не нашёл в них стройности, однородности, формул, характерных для XI~в. Ему удалось выделить лишь 7 статей, которые по своему стилю и синтаксису, могли быть отнесены к этому периоду, однако, французский исследователь не исключал наличия архаизирующей тенденции, которой могли придерживаться легисты-составители «Барселонских обычаев» из окружения Рамона Беренгера IV для того, чтобы придать своей работе более авторитетный характер\footnote{\emph{Bonnassie P. }La Catalogne du milieu du Xe siècle à la fin du XIe siècle. Croissance et mutations d’une société. Toulouse~: Association des publications de l’Université de Toulouse-Le Mirail, 1975. T.~1. P.~31.}. С последним утверждением соглашался такой исследователь-юрист, как С.~Собрекес-и-Видаль, однако, его позиция базировалась на анализе содержания собственно правовых норм\footnote{Цит. по: \emph{Bastardas i Parera J.} Introducció // Usatges de Barcelona. El codi a mitjan segle XII / ed. J. Bastardas i Parerа. Barcelona: Fundació Noguera, 1984 (1991). P.~36.}.

Издатель источника, каталонский историк и палеограф, Дж.~Бастардес-и-Парера, анализируя структуру и композиционный строй источника, пришёл к выводу о наличии в его составе двух преамбул и двух эпилогов, соответственно. Это позволило ему предположить, что текст кодекса, дошедший до нас в рукописной традиции, был результатом компиляции двух памятников. Первым, по его мнению, являлись записи правовых обычаев XI~в. графской курии Рамона Беренгера I, вторым же --- установления «мира и перемирия Божия», утверждённые во время правления Рамона Беренгера III. Создание самой компиляции этот исследователь относил к правлению Рамона Беренгера IV\footnote{\emph{Ibid}. P.~30--31.}.

На сегодняшний день ни одна из вышеуказанных гипотез не была ни доказана, ни опровергнута, поскольку всегда находятся данные, которые могут поставить подобные выкладки под сомнение. Однако стоит отдать должное вышеупомянутым учёным, поскольку появление разных версий о времени возникновения данного источника, стимулировало глубокие лингвистические, исторические, юридические исследования его текста. Отдельно вопроса о датировке текста «Барселонских обычаев» мы касались в статье, выполненной по результатам исследования с использованием квантитативных методов сравнительного анализа текста «Барселонских обычаев»\footnote{\emph{Ясинская (Корнеева) А.Д. }Барселонские обычаи: опыт квантитативного исследования // Проблемы истории и культуры средневекового общества. Материалы XL всероссийской научной конференции студентов, аспирантов и молодых ученых Курбатовские чтения. СПб.: ООО «Скифия-принт», 2021. С.~470--478.}. Данные, полученные в результате контекстного анализа лексического строя «Барселонских обычаев», позволяют говорить о том, что редакция «Барселонских обычаев», опубликованная Ж.~Бастардесом-и-Парерой, сложилась постепенно, но не ранее второй половины XII~в. Аналогичной точки зрения придерживается и большинство современных исследователей.

Среди источников, которые были использованы при составлении «Барселонских обычаев»\footnote{В Главе 1 эти тексты обозначены термином «латинские источники».}, стоит назвать «Вестготскую правду», прямые цитаты из которой мы встречаем в тексте последних\footnote{Usatges de Barcelona. Us. 1 (1--2), 46 (49). P.~48, 86.}. “Lex visigothorum” (более точное название \emph{“}\emph{Liber}\emph{ }\emph{iudiciorum}\emph{”} --- «Книга приговоров») в качестве действующего свода законов, руководства для судей довольно широко применялся в Каталонии в X--XII~вв., поскольку имеются примеры его использования в актах купли-продажи, обмена, дарений и завещаний\emph{ }указанного периода\footnote{\emph{Zimmermann~M. }L'usage du Droit wisigothique en Catalogne du IXe au XIIe siècle... P.~236. См. специальные индексы в этих изданиях грамот: Justícia i resolució de conflictes a la Catalunya medieval. Col·lecció diplomàtica. Segles IX--XI / dir. \emph{J.M.~Salrach} i Marès, T. de Montagut i Estragués. Barcelona: Parlament de Catalunya; Departament de Justícia de la Generalitat de Catalunya, 2018. (Textos Jurídics Catalans; 37. Documents; 2); Justícia i resolució de conflictes a la Catalunya medieval. Col·lecció diplomàtica. Segle XII / dir. \emph{J.M.~Salrach} i Marès, T. de Montagut i Estragués. Barcelona: Parlament de Catalunya; Departament de Justícia de la Generalitat de Catalunya, 2024. (Textos Jurídics Catalans; 38. Documents; 2/3).}. Составители «Барселонских обычаев» несколько раз цитировали «Вестготскую правду» не совсем точно по тексту и в несколько сокращённом виде\footnote{\emph{Попова}\emph{ }\emph{Г}\emph{.}\emph{А}\emph{. }Клятва в законах вестготских королей // Право в средневековом мире. М.:~ИВИ РАН, 2009. С.~219.}. Кроме того, в конце XI--XII~в. на территории Каталонии появились многочисленные переводы «Вестготской правды» на старокаталанский язык\footnote{\emph{Morán i Ocerinjauregui~J. }El proceso de creación del catalán escrito // Aemilianense: revista internacional sobre la génesis y los orígenes históricos de las lenguas romances. San Millán de La Cogolla, 2004. Vol.~1. P.~437.}.

Уже во 2-й статье «Барселонских обычаев» они противопоставлены «Вестготской правде»: \emph{«Когда сеньор Рамон }\emph{Беренгер}\emph{ Старый <…> получил власть над этой землёй и увидел и признал, что во всех делах и тяжбах этой страны }\textbf{\emph{готские законы}}\emph{ не могут соблюдаться <…>, он установил и ввёл ввел обычаи <…>»}\footnote{Usatges de Barcelona. Us. 2 (3): «CVM DOMINVS Raimundus Berengarii uetus <…> habuit honorem et uidit et cognouit quod in omnibus causis et negociis ipsius patrie leges gotice non possent obseruari <...> constituit et misit usaticos...». P.~50.}\emph{. }«Вестготская правда», как минимум, потому что она служила одним из учебных текстов для будущих чиновников, нотариев, писцов, оказала значительное, влияние на правовую культуруXII--XIII~вв.\footnote{\emph{Zimmermann M. }Écrire et lire en Catalogne (IXe--XIIe siècle). Préface de P. Toubert. Madrid, 2003. T.~1. P.~11.}.

Среди других источников, использованных при составлении «Барселонских обычаев», стоит назвать «Составленные Петром извлечения из римских законов». Это своего рода дидактическое сочинение, основанное на выписках и фрагментов из «Дигест» Юстиниана, написанное не ранее середины XII~в. Как установил А.~Гурон, в составе «Составленных Петром извлечений из римских законов» присутствует упоминание “Quoniam egestas” (сокращённой версии «Декрета» Грациана), точное время создания которой -- не ранее 1150~г.\footnote{\emph{Gouron A. }«Petrus» Démasqué.// Revue Historique de Droit Français et Étranger. 2004. Vol.~82. N\textsuperscript{o}. 4. P.~578. URL: www.jstor.org/stable/43851566. (дата обращения: 09.05.2026) В то же время, например, \emph{Д.Ю.~Полдников} отказывается от точной датировки и локализации места составления данного источника: Составленные Петром извлечения из римских законов / пер. с лат. \emph{Д.Ю.~Полдникова}, под ред. \emph{Л.Л.~Кофанова} // Древнее право. Ius antiquum. 2003. №~1 (11). С.~216.}. Это позволяет датировать «Составленные Петром извлечения из римских законов» довольно точно серединой XII~в. Данное руководство к записи городского права было адресовано Одилию, магистрату г.~Валанс (деп. Дром, совр.~Франция). Ещё до определения точной датировки «Составленных Петром извлечений из римских законов», Ю.~Фикер отмечал сходство данного текста и текста «Барселонских обычаев»\footnote{\emph{Ficker J.} Über die Usatici Barchinonae und deren Zusammenhang mit den Exceptiones Legum Romanorum // Mitteilungen des Instituts für Österreichische Geschichtsforschung. 1886. Erg.-Bd.~2. S.\emph{ }242.}.

Другим дополнительным источником стал “Corpus Iuris Civilis”, составленный в VI~в. на территории Византийской империи и объединивший как нормы классического римского права, так и христианские установления, господствовавшие в тот период на территории Восточной Римской империи. Он был заново открыт в г. Равенна (Северная Италии) в ходе подготовки источников для Реформы папы Григория VII в 1070-е гг., примерно в это же время над интерпретацией норм права, сохранившихся в данном кодексе, работали глоссаторы I поколения, Пепо и Иринерий\footnote{\emph{Полдников~Д.Ю. }Договорные теории глоссаторов (XII-XIII~вв.). М., 2008.}. Возможно, что последствием возникновения школ права в Равенне и Болонье стало распространение норм римского права из «Свода~Гражданского права» Юстиниана на территории Прованса, а затем и, возможно, Каталонии. Чтобы проверить данную гипотезу, о которой мы встречаем довольно много упоминаний в научной литературе\footnote{\emph{Ficker J. }Op. cit.}, мы включили полный текст “Corpus Juris Civilis” в число вспомогательных источников.

Ещё одним дополнительным источником для нашего исследования явилась V книга «Этимологий» Исидора Севильского «О законах и временах», посвящённая краткому изложению римского права на основе материала из «Институций» Гая. Известно, что в эпоху высокого Средневековья на территории Пиренейского полуострова в частности, и в Западной Европе в целом известно было около 1000 рукописей данного сочинения\footnote{\emph{Павлов А.А.} «Этимологии или Начала» Исидора Севильского // Адам и Ева. Альманах гендерной истории. М., 2014. №~22. С.~244.}. Учитывая, что «Этимологии» были разделены на книги Брауилионом Сарагосским, епископом Цезареавгусты (совр. Сарагоса), на территории Каталонии в кругу образованного населения они должны были обладать известностью. Это произведение являлось энциклопедией раннесредневекового знания, по свидетельствам исследователей, широко, применявшейся в образовании\footnote{\emph{Уколова В.И. }Рождение средневекового энциклопедизма: Исидор Севильский // Античное наследие и культура раннего средневековья (кон. V -- нач. VII~в.). М. 1989. С.~196--283.}. Изложенные в этом труде представления о римском и каноническом праве, которые, по замечанию Л.Л.~Кофанова, были довольно своеобразными\footnote{\emph{Кофанов Л.Л.} Исидора Гиспальского епископа «Этимологии» или «Начала» // Антология мировой правовой мысли. М., 1999. Т.~2. С.~209.}, также могли повлиять на формирование правовой культуры более позднего периода.

Также в нашей работе при сравнении были использованы тексты четырёх Евангелий. Взять бóльший объём библейского материала не позволили технические возможности\footnote{Подробнее см. Главу 1 настоящего исследования.}. Средневековый человек знакомился с новозаветными текстами довольно рано, в церкви, они часто применялись в обучении грамоте и задавали значимый культурный контекст не только для клириков, но и для мирян. Не последнее место евангельские тексты занимали и в правовой культуре средневекового общества, так И.С.~Филиппов отмечал богатство их правовой лексики, несмотря на их вне- или даже над правовой статус\footnote{\emph{Филиппов И.С.} Библия и римское право: понятия права, закона и справедливости в Вульгате // Философия права Пятикнижия / под ред. \emph{А.А.~Гусейнова}, \emph{Е.Б.~Рашковского}; сост. \emph{П.Д.~Баренбойм}. М.: ЛУМ, 2012. С.~323.}.

Однако возникает вопрос о том, почему в случае с “Liber iudiciorum” каталонские легисты предпочитали явные цитаты, а в отношении других источников сведений о римском праве --- заимствования и неявные аллюзии без ссылок на источник. По нашему мнению, «Вестготская правда» для них оставалась действующим источником права, пусть и постепенно утрачивавшим своё значение. Именно через недостаточность установлений «Вестготской правды» они объясняли саму необходимость составления «Барселонских обычаев»: \emph{«...Ибо граф поступил так на основании Вестготской правды...»}\footnote{Usatges de Barcelona. Us. 2 (3): «...Hoc enim fecit comes auctoritate Libri iudicis qui dicit ...». P.~50.}\emph{.}

В связи с широким распространением и рукописной, и правовой традиции «Вестготской правды» на территории Каталонии в указанный период возникает вопрос о роли «Барселонских обычаев» в правовой культуре региона. Составители довольно чётко отделяли их, как от законов (лат. \emph{--- }\emph{leges}, кат. \emph{--- }\emph{ligs}), которые признаются одним из основных источников права, т. е. «Вестготской правды», так и от норм неписанного права (лат. ---\emph{mores}\emph{, }кат. --- \emph{custums}). Последние, согласно 82(101)-й статье «Барселонских обычаев», могли отличаться на разных территориях: \emph{«…согласно законам или согласно обычаям их земли»}\footnote{Ibid. Us. 82 (101): «...secundum leges uel secundum istius patrie mores». P.~120.}.

Термин \emph{“}\emph{lex}\emph{” }(лат.---закон) в «Барселонских обычаях» используется в двух основных значения, с одной стороны, это, несомненно, \emph{“}\emph{Lex}\emph{ }\emph{visigothorum}\emph{”}, к тексту которого, как мы писали выше, обращались неоднократно при составлении документов, в спорных случаях уголовного законодательства и т.д., с другой же, это обозначение религиозной принадлежности, как, например, в 70 статье\footnote{Ibid. Us. 70 (75). P.~108.}, что довольно типично для источников данного региона.

В то же время в 123-й статье «Барселонских обычаев» мы встречаем следующую иерархию источников права: \emph{«…Установили выше названные принцепсы, чтобы судили, согласно «Обычаям», где «Обычаев» будет недостаточно, возвращались к законам и к решениям принцепсов, судей или курии»}\footnote{Ibid. Us. 123 (81): «…constituerunt prelibati principes secundum usaticum esse iudicata, et, ubi non sufficerent usatici, reuerterentur ad leges et ad principis arbitrium, eiusdemque iudicium atque curie». P.~158.}.

Таким образом, правовая культура, в которой существовали «Барселонские обычаи» была само по себе многослойной и «многоголосной». Составление этого свода, по словам своих составителей, было призвано дополнить уже существующие нормы в тех областях права, которые были слабо разработаны в «Книге приговоров», например, в области феодальных отношений: »\emph{«…потому что в законах не находим оммажа [не говорится об }\emph{оммаже.---}\emph{А.K.]»}\footnote{Ibid. Us. 123 (81):  «<…> quia in legibus non inuenitur hominiaticum <…>». P.}. Это не было единственной целью составления данного законодательного кодекса. Общеизвестно, что любая фиксация обычного права, существующего в устной форме, предполагает его стандартизацию, хотя бы в рамках записи текста для определённой территории, языковых особенностей и т.д.

В этом можно увидеть некий политический проект и правовое обоснование т.н. теории принципата в Каталонии\footnote{Например, \emph{Kagay D.J. }Princeps Namque: Defense of the Crown and the Birth of the Catalan State // Mediterranean Studies. Pennsylvania, 1999. Vol.~8. P.~55--87; \emph{Kosto A.J. }The limited impact of the “Usatges de Barcelona” in Twelfth-century Catalonia // Traditio. Cambridge: Cambridge, 2001. Vol.~56. P.~53--88.}. Именно в XII~в. при Рамоне Беренгере III и Рамоне Беренгере IV власть графов-королей значительно укрепилась в Каталонии. Ни «Вестготская правда», ни местные правовые установления в их изначальном виде не могли служить фундаментом для построения политической системы, в то время как «Барселонские обычаи», составление которых официально приписывалось курии Рамона Беренгера I, подходили для такой роли как нельзя лучше. При этом, по нашему мнению, стоит избегать переоценки значения политического фактора составления «Барселонских обычаев».

К XIII~в. в канцеляриях Каталонии всё большее число документов составлялись на старокаталанском языке, который к тому времени ещё не был стандартизирован, но весьма активно развивался. Именно к этому периоду, по мнению большинства исследователей, в том числе и издателя «Барселонских обычаев», Х.~Бастардеса-и-Пареры, и относятся многочисленные переводы этого свода обычного права на старокаталанский язык\footnote{\emph{Bastardas i Parera J.} Introducció... . P.~9--10.}. Со временем число таких переводов росло. Вероятно, на территории Принципата Каталонии в XIII--XIV~вв. существовало несколько переводческих центров. Именно это и обусловило появление 18 дошедших до нас рукописей старокаталанской версии «Барселонских обычаев», которые тоже довольно сильно различаются между собой\footnote{\emph{Diez de Revenga P., García Diaz I. }Op.cit. P.~82.}: некоторые из них придерживаются архаизированного порядка слов, приближенного к латинскому оригиналу\footnote{\emph{Meyer-Hermann R.} Auxiliar-Konstruktionen im llatí medieval katalanischer Urkunden (10.--13.~Jh.) und in den Usatici Barchinonae / Usatges de Barcelona (12.--15. Jh.). // Zeitschrift für romanische Philologie. 2019. 135. S. 507--534.}, другие же отличаются большим числом инноваций и заимствований\footnote{\emph{Diez de Revenga P., García Diaz I.} Op.cit. P.~83.}. В этот же период сложилась и традиция глосс и комментариев к тексту «Барселонских обычаев», анализ которой выходит за рамки нашего исследования.

Текст «Барселонских обычаев» разделён на статьи, в каждой из которых содержится одна или несколько правовых норм, обычно, но не всегда, связанных тематически. Порядок взаимного расположения рубрик и их нумерация весьма существенно отличаются в различных рукописях источника и, соответственно, в текстах изданий\footnote{\emph{Mor~C.G}. En torno a la formación del texto de los "Usatici Barchinonae" // Anuario de historia del derecho español. Nº 27--28. Madrid, 1957--1958. P.~417--423.}. Взаимное расположение статей «Барселонских обычаев» от рукописи к рукописи не было устойчивым, поэтому нумерация, используемая Х.~Бастардесом-и-Парерой\footnote{\emph{Bastardas i Parera J.} Introducció. // Usatges de Barcelona. Barcelona, 1991. P.~7--38.}, весьма условна и включает два номера: первый по латинской рукописи P, второй по старокаталанской рукописи~K.А.~Иглесиаc~Ферейрос выступает категорически против такого деления, утверждая, что из-за совпадений в статьях «преамбулы», статьях~№~140, 188, 190 и далее в поздних рукописях, их можно датировать более поздним периодом, в частности, правлением Жауме I Завоевателя\footnote{\emph{Iglesia }\emph{Ferreirós}\emph{ A.} The birth of the Usatici // Imago Temporis --- Medium Aevum. 2011. Vol.~5. P.~132.}.

До наших дней сохранилась 61 рукопись\footnote{\emph{Pérez Martín~A. }Hacia una edición crítica del texto latino de los Usatges de Barcelona. // Glossae: European Journal of Legal History. Nº 7. 1995. P.~9--32.}, содержащая текст «Барселонских обычаев»: 43 манускрипта на латинском языке (35~самостоятельных рукописей, 8 рукописей-«цитирований», помещенных в составе трактатов по праву и компиляций) и 18 --- на каталанском (12~самостоятельных рукописей, 6~рукописей-«цитирований» в составе трактатов по праву и компиляций); что для средневекового текста довольно много (например, примечательно, что это число превышает количество дошедших до нас рукописей \emph{“}\emph{Liber}\emph{ }\emph{iudiciorum}\emph{”}\footnote{\emph{Попова~Г.А.} «Вестготская правда» («Книга приговоров»): рукописная традиция и история издания // Вестготская правда (книга приговоров). Латинский текст. Перевод. Исследование / Oтв. ред. \emph{О.В.~Ауров}, \emph{А.В.~Марей}. М., 2012. С.~131.}!).

Первые рукописи этого памятника на старокаталанском языке относятся, по палеографическим данным, к XIII~в. (при том, что древнейший латинский манускрипт датируется концом XII~в.). Однако, по мнению А.~Переса Мартина, выяснение отношений родства между известными рукописями «Барселонских обычаев» и построение стеммы, раскрывающей характер этих связей, требуют отдельного тщательного исследования\footnote{\emph{Pérez Martín~A. }Hacia una edición crítica del texto latino de los Usatges de Barcelona. // Glossae: European Journal of Legal History. Nº 7. 1995. P.~29.}.

Сохранению наследия рукописной традиции «Барселонских обычаев» способствовала децентрализованная система архивов на территории сначала Арагонской короны, а потом и Испанского королевства. Самым древним из манускриптов считается относимый палеографами к 1190-м гг. MS. Lat.4792\footnote{BNF. Départament des Manuscripts. MS. Lat. 4792. v2.}, хранящийся в Отделе рукописей Национальной Библиотеки Франции.

Тексты многих рукописей были использованы в публикациях источника, первая из которых была осуществлена в 1495~г.\footnote{Usatges de Barcelona e Constitucions de Catalunya. Barcelona. Pere Miquel y Diego de Gumiel.1495.}. «Барселонские обычаи неоднократно переиздавали: в XIX---Ш.~Жиро\footnote{\emph{Giraud Ch. }Op.cit.}, в XX~в.--- 6 раз: в 1907, 1913, 1933, 1984, 1991~гг.\footnote{Gudiol J. «Usatges de Barcelona» // Anuari de l’Institut d’Estudis catalans. №~1. 1907. P.~285--334; \emph{d’Abadal i Vinyals R., Valls i Taberner F.} Usatges de Barcelona. Textes de Dret Català. Barcelona, 1913; \emph{Rovira i Ermengol J. }Usatges de Barcelona i commemoracions de Pere Albert. Barcelone, 1933. P; Los Usatges de Barcelona: Estudios, comentarios,y edición bilingüe del texto. Malaga: Universidad de Malaga, 1984; Usatges de Barcelona. El codi a mitjan segle XII. Editado por J. Bastardas i Parerа. Barcelona., 1984 (1991).}, в XXI~в. --- в составе сборника «Исторические законы Каталонии». В конце XX---начале XXI~вв. эта традиция поддерживалась многочисленными комментированными публикациями отдельных рукописей в журнале \emph{“}\emph{Initium}\emph{. }\emph{Revista}\emph{ }\emph{catalana}\emph{ }\emph{d}\emph{'}\emph{hist}\emph{ó}\emph{ria}\emph{ }\emph{del}\emph{ }\emph{dret}\emph{”}\footnote{\emph{Iglesia Ferreirós A.} El manuscrito latino 4792 de la Biblioteca Nacional de París. Usatges y Liber Iudiciorum // Initium: Revista catalana d’historia del dret. 2000. №~5. P.~643--826;\emph{ Idem.} Introducción a una edición ideal de Usatici y Glosas // Initium: Revista catalana d’historia del dret. 2000. №~14. P.~3--194; \emph{Idem.} De usaticis quomodo inventi fuerunt // Initium: Revista catalana d’historia del dret. 2001. №~6. P.~25--212;\emph{ Idem. }Los Usatges de Barcelona: una nota crítica // Initium: Revista catalana d’historia del dret. 2001. №~6. P.~385--408; \emph{Idem}. Usatici, paces, constitutiones. Edición del ms. lat. 4249 de la Bibliothèque Nationale de Paris // Initium: Revista catalana d’historia del dret. 2001. №~6. P.~481--714; \emph{Idem. }Glosas y Usatges el ms. BNP Latin 4670 A Edición // Initium: Revista catalana d’historia del dret. 2002. №~7. P.~363--847; \emph{Idem. }Biblioteca Apostólica Vaticana Ms. Ottob. Lat. 3058 de los usatges glosados. Edición // Initium: Revista catalana d’historia del dret. 2003. №~8. P.~511--894 и др.}\emph{ }под руководством А.~Иглесии~Ферейроса и его исследовательской группы. Однако полное комментированное издание «Барселонских обычаев», о необходимости которого неоднократно пишут А.~Иглесиас~Ферейрос\footnote{\emph{Idem}. Introducción a una edición ideal de Usatici y Glosas // Initium: Revista catalana d’historia del dret. 2009. №~14. P.~3--194.} и другие исследователи, даже по результатам такой тщательной и кропотливой работы всё ещё не увидело свет.

Помимо изданий текста на латинском и каталанском языках, в XIX---XX~вв. было осуществлено два перевода «Барселонских обычаев» на европейские языки: кастильский\footnote{\emph{Nolasco Vives y Cebria P. }Traducción al Castellano de los Usajes y demas derechos de Cataluña. 4 vols. Barcelona. 1861.} и английский\footnote{The Usatges of Barcelona: the fundamental law of Catalonia / trans. \emph{D.J.~Kagay}. Philadelphia: University of Pennsylvania Press, 1994.}. Второй представляет для нас особый интерес, поскольку снабжён подробной вступительной статьёй, приложениями и детальными комментариями Д.~Дж.~Кэгэя\footnote{\emph{ }\emph{Kagay D.J}\emph{.} Introduction and Notes // The Usatges of Barcelona: The Fundamental Law of Catalonia. Philadelphia: University of Pennsylvania Press, 1994. P.~1--60.}, однако, это издание содержит отдельные неточности и требует дополнений и пояснений, которые зачастую обусловлены тем, что автор практически не обращается к старокаталанскому тексту памятника. Перевода текста «Барселонских обычаев» на русский язык не существует, однако только отдельные установления были переведены В.К.~Пискорским, Л.Т.~Мильской и другими исследователями\footnote{\emph{Пискорский}\emph{ В.К.} Крепостное право в Каталонии в средние века. М., 2019.}.

В своей работе мы пользовались публикацией текста источника Ж.~Бастардеса-и-Пареры, в которой, издатель на основе сопоставления данных 5 наиболее ранних манускриптов и первого печатного издания «Барселонских обычаев» осуществляет реконструкцию латинского текста «Барселонских обычаев»\footnote{\emph{Bastardas i Parera J. }Introducció // Usatges de Barcelona. El codi a mitjan segle XII / ed.~J.~Bastardas i Parerа. Barcelona: Fundació Noguera, 1984 [1991]. P.~25.}. Она отличается и тем, что параллельно этому тексту приводится восстановленный текст отдельных статей на старокаталанском языке по рукописи начала~XIV~в. из Архива Арагонской Короны\footnote{ACA. Cancillería. Papeles por incorporar. Legislació. Caixa 1. Núm.1.}. Ж.~Бастардес-и-Парера отмечает, что это копия более раннего манускрипта, который, по лингвистическим данным, может быть датирован XIII в\footnote{\emph{Bastardas i Parera J. }Introducció // Usatges de Barcelona. El codi a mitjan segle XII / ed.~J.~Bastardas i Parerа. Barcelona: Fundació Noguera, 1984 [1991]. P.~23.}. Отдельно издатель отмечал существование нескольких переводов «Барселонских обычаев» на старокаталанский язык, вероятнее всего, выполненный в разных местах, однако для публикации он выбрал именно эту версию ввиду того, что язык в ней наиболее соответствует литературным нормам своего времени, а порядок расположения статей сопоставим с тем, что фиксируется в большинстве латинских рукописей\footnote{\emph{ Ibid.} P.~24.}. Это издание снабжено наиболее проработанным научным аппаратом: обширным введением, обозначениями разночтений в различных рукописях, комментариями издателя на современном каталанском языке к наиболее сложным для чтения местам.

Статьи латинской версии «Барселонских обычаев» по изданию Х.~Бастардеса-и-Пареры в 2020~г. были опубликованы в составе \emph{“Corpus }\emph{Documentale}\emph{ }\emph{Latinum}\emph{ }\emph{Cataloniae}\emph{”}\footnote{Corpus Documentale Latinum Cataloniae. URL~: \url{https://gmlc.imf.csic.es/glossarium/codolcat} , дата обращения: 03.04.2026.}. Однако данная онлайн-публикация в большей степени ориентирована на актовый материал, что не позволяет обеспечить выдачу полного текста статьи «Обычаев», ограничивая объём выдачи определённым числом знаков перед и после искомого слова или набора слов.

Что касается вопроса об авторстве «Барселонских обычаев», то он был поставлен ещё в XIX~в. и также всерьёз обсуждается в научной литературе. На наш взгляд, говорить о какой-либо конкретной персоналии в качестве автора данного правового свода нам кажется преждевременно до того момента, пока окончательно не будет решён вопрос с датировкой самого источника. И сегодня мы можем лишь сказать, что автор входил в Совет при дворе графа и принадлежал к одной из знатных семей в окружении графов Барселонских. Впрочем, установлено, что образование этот человек получил на территории Каталонии, Прованса или Северной Италии, об этом нам говорит, например, не совсем точное, скорее всего, по памяти, цитирование фрагментов из «Вестготской правды», использование наряду с традиционными римскими терминами лексики каталонского и окситанского происхождения.

Некоторыми исследователями предпринимались попытки точно определить его личность: одни называли среди возможных авторов первой редакции Понса Бонфиля Марка, куриального судью и легиста времён правления Рамона Беренгера I\footnote{\emph{Zimmermann M. }La représentation de la noblesse dans la version primitive des Usatges de Barcelone (milieu du XIIe siècle) // Cahiers de linguistique et de civilisation hispaniques médiévales. N°25. 2002. P.~15.}, другие считали, что в создании «Барселонских обычаев» мог принимать участие лично граф Рамон Беренгер~III или кто-либо из приглашённых им ко двору провансальских легистов\footnote{\emph{Sobrequés i Vidal S. }Història de la producció del dret català fins al decret de nova planta. Girona, 1978. P.~21.}. Однако, как нам кажется, подобные гипотезы не могут быть никоим образом верифицированы, а это обстоятельство свидетельствует о том, что текст «Барселонских обычаев» требует дополнительного исследования. Более целесообразно в данном случае говорить о коллективном авторстве, о том, что в составлении «Барселонских обычаев» принимали участие принадлежавшие к разным поколениям легистов и придворных. Более того, «авторами» данного источника отчасти были безымянные каталонцы, некогда пользовавшиеся собственно обычаями, которые затем были переработаны и записаны в графской курии.

Структура дошедшего до нас текста «Барселонских обычаев» довольно сложна: в ней выделяются несколько пластов, различных по времени возникновения. Традиционно исследователи выделяют текст «архетипа», состоящий из первых 138 статей, записанных, по мнению Х.~Бастардеса-и-Пареры, не позднее 1150~г., они содержатся во всех рукописях памятника\footnote{\emph{Bastardas i Parera J.} Introducció ... P.~9--10.}. «Официальная редакция» памятника в XIV--XV~вв., закреплённая в \emph{“}\emph{Constituciones}\emph{ }\emph{de}\emph{ }\emph{Cataly}\emph{ñ}\emph{a}\emph{”}, включает уже 173--174 статьи (новые законодательные постановления добавлялись к уже существующему тексту «Обычаев» постепенно, например, при Жауме I Завоевателе)\footnote{\emph{Gonzalvo i Bou G.} Els jueus i els Usatges de Barcelona // Barcelona Quaderns D'Història. 1996. №~2. P.~118.}. А.~Иглесиас~Ферейрос насчитывает всего 304 статьи, когда-либо входивших в текст «Барселонских обычаев», что говорит, что у текста «Обычаев» были другие источники расширения: \emph{глоссы и комментарии легистов}, фрагменты которых со временем стали восприниматься как часть самого правового памятника\footnote{\emph{Iglesia }\emph{Ferreirós}\emph{ A. }The birth of the Usatici // Imago Temporis --- Medium Aevum. 2011. Vol.~5. P.~120--121.}. В нашем исследовании, мы будем опираться на текст «Обычаев», включающий 138 статей, приведённый в издании Х.~Бастардеса-и-Пареры.

Исследование «Барселонских обычаев» на двух языках в наиболее ранних редакциях отличается высокой репрезентативностью, поскольку именно анализ перевода позволяет увидеть правовую культуру в момент рефлексии её носителей: каким образом тот или иной легист-переводчик трактовал латинские термины, что он опускал, а что, наоборот, считал наиболее важным. Однако, надо понимать, что мы имеем дело лишь с фиксацией процесса перевода и интерпретации в его начальной и конечной точках, что, несомненно влечёт за собой информационные потери: от нас скрыта последовательности логических операций, происходившая в сознании средневекового человека и приводившая его к тому или иному переводческому решению. Следует избегать также и обобщений, поскольку в данном случае мы имеем дело c творческим выражением языковой личности переводчика, какие-то из принятых им переводческих решений были субъективны или отражали его личное понимание тех или иных терминов. Однако тот факт, что он их использует, уже является для нас показательным и в какой-то мере, с поправкой на регулярность определённого рода изменений и сравнение различных переводов, мы можем говорить о крайне высокой репрезентативности «Барселонских обычаев» в отношении интересующей нас проблематики.

\subsection*{§2. «Обычаи Льейды» (1228)}
\addcontentsline{toc}{subsection}{§2. «Обычаи Льейды» (1228)}

Составление «Обычаев Льейды» (Consuetudines ilerdenses) традиционно датировали 1228~г.\footnote{Ф.~Вальс-и-Табернер считал, что традицию датировки составления «Обычаев Льейды» 1228~г. задал испанский историк Вильянуэва. \emph{Valls y Taberner F.} Las «Consuetudines ilerdenses» (1227) y su autor Guillermo Botet. Barcelona: [s.n.], 1913. P.~14.} Однако стоит заметить, что Ф.~Вальс-и-Табернер предложил датировать составление этого свода годом ранее по следующей причине: Гульельм Ботет указал, что во время составления «Обычаев Льейды» он являлся городским консулом. В документе от 16 ноября 1226~г. Г.~Ботет фигурирует в качестве консула, что и в предисловии к «Обычаям Льейды». Известно, что городской консул избирался ежегодно, а занимать должность несколько лет подряд было нельзя. Поэтому Г.~Ботет не мог быть консулом в 1228~г. Также Ф.~Вальс-и-Табернер нашёл рукопись «Обычаев Льейды», датированную 1227~г.\footnote{\emph{Valls y Taberner F. }Op. cit. P.~14--15.}. В одной из сохранившихся рукописей и одной, недошедшей до нас присутствует датировка составления «Обычаев Льейды» 1232~г., однако, и Ф.~Вальс-и-Табернер, и М.Б.~Пилар-Лоссерталес-и-де-Вальдеавейяно, посчитали её ошибкой писца\footnote{Цит. по: \emph{Barrero García A.M.} Las Costumbres de Lérida, Horta y Mirabet // Anuario de historia del derecho español. 1974. №~44. P.~487.}.

Необходимо учитывать, что так или иначе речь идёт о фиксации уже сложившейся к этому времени правовой практики города и его окрестностей. В тексте «Обычаев Льейды» получил отражение более ранний пласт норм, восходящий к периоду после завоевания Льейды в 1149~г., и статьи, по языку соотносящиеся уже с правовой культурой первой трети XIII~в.

Авторство компиляции «Обычаев Льейды{\unicodefallback ̱}» практически все исследователи вслед за оригинальным текстом приписывают городскому консулу Гульельму Ботету, который попал под влияние университетской римско-канонической традиции. Он был тесно связан с городским советом и судом, служил посредником между устной правовой традицией общины г. Лериды и её окрестностей и требованиями королевской канцелярии.

Что касается источников формирования норм, внешняя критика выделяет несколько ключевых пластов:

\begin{itemize}
\item во-первых, это поселенная хартия Льейды 1149~г., пожалованная вскоре после христианского завоевания города, в которой определялись статус поселенцев, основы муниципального управления, распределение земель, фискальные и военные обязательства населения;
\item во-вторых, королевские привилегии XII--XIII~вв., дарованные Льейде и её округе графами Барселонскими и арагонскими, затрагивавшие компетенции городских органов власти, процессуальное право, налоговые льготы и статус отдельных групп населения.
\item в-третьих, важнейшим источником считались и локальные обычаи, сложившиеся в городской среде и судебной практике: типичные способы отчуждения имущества, устойчивые правила кредитования и залога, наследования, санкции за насилие, кражу и оскорбления. В тексте «Обычаев Льейды» встречаются формулы и конструкции, явно навеянные римско-каноническим правом и текстом «Барселонских обычаев», что свидетельствует о влиянии учёной традиции на запись локального обычая.
\end{itemize}
Отдельно исследователи вслед за текстом преамбулы указывают на значение «Барселонских обычаев» XII~в. как одного из ориентиров для городского права Льейды\footnote{\emph{Oliver y Esteller B.} Historia del derecho en Cataluña, Mallorca y Valencia. Código de las costumbres de Tortosa. Madrid: Impr. de M. Ginesta, 1876. T.~1. P.~981.}, кроме того, как дополнительные источники права «Обычаи Льейды» признавали «Вестготскую правду» и римское право\footnote{\emph{Valls y Taberner F.} Op. cit. P.~22.}. Барселона к этому времени уже имела развитую традицию, частично закреплённую в письменных формах, и её правовые своды служли моделью для других городов каталоно-арагонского ареала. Сравнение «Обычаев Льейды» с барселонскими обычаями выявляло параллели в регулировании городского управления, имущественных и обязательственных отношений, а также процессуальных правил.

Первоначально «Обычаи Льейды» были записаны на латинском языке с использованием отдельных старокаталанских терминов, это свидетельствует о том, что составители кодекса стремились в подражание «Барселонским обычаям» придать кодексу большую легитимность и авторитет.

Рукописная традиция «Обычаев Льейды» восходит к позднесредневековым спискам XIV--XV~вв.; автографа или раннего оригинала не сохранилось. Ф.~Вальс-и-Табернер указывал на существование 4 рукописей «Обычаев Льейды» (А\footnote{Archivo municipal de Lérida. Vol.~1376 (antes 1372). Libro verde menor. Ff. 69--92.}, B\footnote{Archivo dela Catedral de Lérida. Armario L. Cajón M. Usatici et Consuetudines Cathaloniete Constitutiones. Provincial et sinodales,. Ff. 123--133.}, C\footnote{Biblioteca de la familia de Dalmases, en Barcelona. №~50. Consuetudines Ilerdenses. cum glosulis. Fori Aragonis tam antiqui quam novi, cum suis glosis (2). Ff. 1--11.} и D\footnote{Biblicteca Nacional de Madrid. №~865 (antes C.~193). Ff. 1--22.}), ни одна из которых не была ни оригиналом, ни его прямой копией. Кроме того, он указал ещё 5 рукописей, которые ему не удалось обнаружить в архивах, но которые упоминались в описях имуществ различных лиц\footnote{\emph{Valls y Taberner F.} Op. cit. P.~29.}. Особый интерес в этом контексте представляет то обстоятельство, что две из них по сведениям Ф.~Вальса-и-Табернера представляли собой перевод «Обычаев Льейды» на старокаталанский язык\footnote{Ibid.}. Таким образом, старокаталанская версия этого правового компендиума до нас не дошла.

В свою очередь к 1946~г. М.Б.~Пилар-Лоссерталес-и-де-Вальдеавейяно указала на ещё одну рукопись XV~в. (назовём её E\footnote{La Biblioteca Central de la Diputación de Barcelona. Ms. 702. Ff. CV--CXVIII.}). На ранних этапах текст бытовал преимущественно на латинском языке, но по мере переписей и переводов усиливается присутствие романских (каталанских) элементов. Для внешней критики это означает необходимость осторожной реконструкции первоначального слоя текста 1227--1228~гг., учёт возможных поздних интерполяций и адаптаций норм к изменившейся практике.

На момент момент создания диссертационного исследования Ф.~Вальса-Табернера было известно 2 издания «Обычаев Льейды» в составе крупных коллекций исторических источников: П.~Виллануэвы\footnote{\emph{Villanueva J., Villanueva J.L. }Viage literario á iglesias de España~: Le Publica con algunas observaciones. Madrid~: Imprenta de la Real Academia de la Historia, 1803. T. XVI. P.~160--196.} и Х.~Плеяна-де-Порты\footnote{\emph{Pleyan de Porta J. }Apuntes de historia de Lérida~: ó sea compendiosa reseña de sus mas principales hechos desde la fundación de la ciudad hasta muestros tiempos. Lérida: Carruez, 1873. P.~495--528.}. Среди недостатков этих изданий Ф. Вальс-Табернер выделил опору лишь на рукопись Е без привлечения иных и некритический характер. Эта проблема была решена критическим изданием «Обычаев Льейды», подготовленным испанской исследовательницей Б.М.~Пилар-Лоссерталес-и-де-Вальдеавейяно в 1946 году по материалам всех известных рукописей\footnote{По возможности мы обращались именно к этому изданию, однако, оно было доступно нам лишь частично: Costumbres de Lérida / ed. \emph{M.B.~Pilar}~Loscertales de Valdeavellano. Barcelona: Universidad de Barcelona, Facultad de Derecho, 1946. P.~10--12.}. В~1997~г. «Обычаи Льейды» были переизданы Дж.Дж.~Бускетой-и-Риу и Е.~Гонсалес с обширным предисловием, латинской транскрипцией и факсимильным изданием рукописи из городского Архива Льейды, также в этом издании был дан перевод текста правового памятника на современный каталанский язык\footnote{К~сожалению, это издание осталось для нас недоступным: Els costumns de Lleida / ed. por \emph{J.J.~Busqueta} Riu, E. González. Lleida: Ajuntament de Lleida, 1997.}.

В контексте настоящего диссертационного исследования важно упомянуть и рецепцию «Обычаев Льейды» в городских установлениях Орты, Миравета и ряда других каталонских коммун XIII~в., находившихся под управлением ордена тамплиеров. Это свидетельствовало о превращении свода в образец муниципального обычного права. Тексты местных обычаев фиксируют как прямые заимствования отдельных норм и формулировок из «Обычаев Льейды», их адаптацию к специфике социально-экономической структуры конкретных городов. При этом «Обычаи Льейды» служили не только источником конкретных правовых решений, но и образцом организации городского правового пространства, задавая общие рамки для регулирования статуса горожан, имущественных отношений и судебной практики. Сравнительный анализ Обычаев Орты, Миравета и Льейды позволил А.М.~Гарсии-и-Барреро и Ж.Серрано-Дауре выделить механизмы трансформации исходного текста в процессе переноса, а также степень участия королевской власти в стандартизации муниципального обычного права Каталонии\footnote{\emph{Barrero García A.M. }Op. cit. P.~494--498; \emph{Serrano Daura J. }La projecció del dret de Lleida // Revista de Dret Històric Català. 2004. №~3. P.~233--258. «Обычаи Орты», «Обычаи Миравета» и др. мы рассмотрим далее.}.

Этот сборник позволяет проследить, каким образом местные правовые традиции интегрировались складывающейся каталонской правовой системой и взаимодействовали с королевским законодательством и римско-каноническими моделями. Обращение к тексту «Обычаев Льейды» даёт возможность реконструировать регулирования социальных конфликтов, статус различных групп населения и специфику правосознания горожан и легистов. В контексте диссертации данный источник служит для анализа соотношения формализованных правовых норм и живой правовой практики, а также для выявления региональных особенностей правовой культуры Каталонии в рассматриваемый период.

\subsubsection*{§2.1. Семейство «Обычаев Льейды»}
\addcontentsline{toc}{subsubsection}{§2.1. Семейство «Обычаев Льейды»}

\paragraph*{§2.1.1. «Обычаи Тарреги» (1242)}

«Обычаи Тарреги»\footnote{\emph{Font Rius J.M. }Costumbres de Tárrega // Anuario de Historia del Derecho Español. 1953. T.~23. P.~429--444.} являлись одним из самых лаконичных сводов местного обычного права: в состав этого памятника вошли всего 25 статей. Отчасти это объясняется его характером, на что указывал издатель этого свода обычного права, Дж.М.~Фонт-Руис, отмечая, что по формальным признакам можно отнести этот документ к жанру королевской привилегии, пожалованной населению города. При этом содержание правовых норм «Обычаев Тарреги» указывает на их активное использование на территории этого поселения.

В период, предшествующий составлению обычаев, Таррега была небольшим замком, отвоёванным у мавров Рамоном Беренгером I в XI~в., затем до 1219~г. упоминания об этом поселении имели спорадический характер и были связаны с хозяйственной деятельностью\footnote{\emph{Serrano Daura J.} La projecció del dret de Lleida // Revista de Dret Històric Català. 2004. №~3. P.~233--258.}. Именно рост экономической активности и усиление городского совета стали предпосылками к письменной фиксации местных обычаев.

«Обычаи Тарреги» были пожалованы городу Жауме I Завоевателем 8~марта 1242~г., как полагает Дж.М.~Фонт-Руис, этому предшествовало ходатайство местной администрации. При составлении заимствовались нормы из более авторитетных правовых памятников: «Обычаев Лериды» и «Барселонских обычаев».

Известно 4 рукописи этого источника, обозначаемые следующей номенклатурой:

А -- оригинал, хранящийся в Библиотеке Каталонии (XIII~в.)\footnote{Biblioteca Central de la Deputación de Barcelona. Sección de Archivo. Pergarnino 04458. Caja. 4-VI-6.};

B -- копия в Историческом архиве г. Тарреги в составе «Книги Привилегий Тарреги» (XIII~в.)\footnote{Archivo Municipal de la villa de Tarrega. Libro de Privilegios. Vol.~1. Fol. 1\textsuperscript{o}};

C -- ещё одна копия в составе «Книги Привилегий Тарреги» (XIII~в.)\footnote{Archivo Municipal de la villa de Tarrega. Libro de Privilegios. Vol.~2. Fol. 1\textsuperscript{o}};

D -- копия 1357~г. в Архиве капитула г. Сео-де-Уржель\footnote{Biblioteca de Cataluña. Biblioteca Capitular de Urgel.Núm. 143. Fol. 129r--133v.}

Также «Обычаи Тарреги» до публикации Дж.М.~Фонт-Руисом в 1953~г. (этим изданием мы пользовались в настоящей работе) были изданы в усечённом варианте Л.~Сарретом-и-Понсом в 1930~г. в переводе на каталанский язык\footnote{Privilegis de Tàrrega / [a cura de] Lluís Sarret i Pons, Pvre ; A. Duran i Sanpere (Pr.) Tarrega : Impr. F. Camps Calvet, 1930, P.~11--15.}.

Текст «Обычаев Тарреги» записан на латинском языке с минимальным количеством каталанских терминов, кроме того, он практически не структурирован тематически в ранних рукописях, при этом выделялись преамбула, основная часть, разделённая на статьи, и заключительная.

Согласно 24-й статье, «Обычаи Тарреги» занимали место дополнительного источника права по отношению к «Барселонским обычаям»,вестготским законам и римскому праву\footnote{\emph{Font Rius J.M. }Costumbres de Tárrega // Anuario de Historia del Derecho Español. 1953. T.~23. P.~443.}. В то же время, как полагал Дж.М.~Фонт-Руис, это положение могло быть заимствовано из «Обычаев Льейды» и являлось довольно характерным для городских центров этого региона, вновь заселённых преимущественно выходцами из центральных районов Каталонии, где «Вестготская правда», а позднее и «Барселонские обычаи», воспринимались как авторитетные источники права\footnote{Ibid. P.~443.}.

Таким образом, «Обычаи Тарреги» --- один из источников для изучения становления правовой культуры Каталонии XIII~в., процессов рецепции римского права и формирования городского самоуправления на стыке традиции и внешних правовых влияний.

\paragraph*{§2.1.2. «Обычаи Орты» (1296)}

Следующим памятником, связанным с правовой традицией «Обычаев Льейды», являются «Обычаи Орты», утверждённые 16 апреля 1296~г. в замке Орты, ныне Орта-де-Сан-Жоан. Источник возник в пределах владений ордена тамплиеров на юго-западной окраине Каталонии, в районе Терра-Альта. В данном случае основной причиной создания местной записи обычного права стал спор между местным представителем ордена тамплиеров и жителями Орты о том, какие именно обычаи должны применяться в Орте и её округе\footnote{\emph{Serrano Daura J. }La projecció del dret de Lleida // Revista de Dret Històric Català. 2004. Т.~3. P.~246.}. Сам памятник прямо характеризуется как соглашение между магистром ордена тамплиеров Беренгером де Кардона и общиной Орты. Оно было призвано заменить ранее действовавшие «Обычаи Льейды» и одновременно оговорить невозможность прямого применения общего права во всех делах без ущерба для вассалов Ордена.

Памятник состоит из 81 статьи при этом деление на книги и иные более крупные фрагменты отсутствует, не наблюдается также тематически обусловленной структуры расположения статей. В тексте «Обычаев Орты» соседствуют нормы о правовом положении жителей, пользовании пастбищами, водами, лесами и дорогами, рыночной торговле, мерах и весах, долговых отношениях, залоге, судебном процессе, уголовных наказаниях, наследовании, брачно-имущественных отношениях и компетенции местных должностных лиц. Такое построение характерно для локальных правовых сводов XIII~в.: составители текста не стремились к отвлечённой классификации норм по отраслям права, и при этом старались зафиксировать практически значимые ситуации, возникавшие в повседневной жизни общины. При этом порядок изложения не следует считать признаком «архаичности» источника: напротив, многие процессуальные положения демонстрируют знакомство его составителей с другими сводами обычного права.

До нас дошёл латинский оригинал, и именно латинская редакция является наиболее ранней. Вместе с тем Х.~Серрано-Даура указывает на возможное существования старокаталанского текста, однако, не раскрывает и не конкретизирует свои доводы\footnote{\emph{ Idem. }Op. cit... P.~247.}. По этой причине говорить о старокаталанской версии памятника не представляется возможным. Корректно охарактеризовать данный памятник как латинский текст, возникший в каталонской правовой среде и, возможно, связанный с несохранившейся старокаталанской версией.

До наших дней сохранился оригинал этой свода, который хранится в Национальном историческом архиве Испании (Мадрид), в фонде энкомьенды Орты\footnote{AHN. Órdenes Militares. San Juan de Jerusalén. Encomienda de Horta. Carp. 672. Doc. 8.}; известны также позднейшие копии 1570 и 1575~гг.\footnote{AHN. Castellanía de Amposta. Caja 8125. Doc. 8; AHN. Órdenes Militares. San Juan de Jerusalén. Encomienda de Horta. Carp. 672. Doc. 9.}. Памятник был впервые опубликован Дж.~Котс-и-Горксом в 1930~г. по тексту оригинальной грамоты, обнаруженному издателем в Национальном историческом архиве Мадрида\footnote{Cots i Gorchs J. Les consuetuds d’Horta (avui Horta de Sant Joan), a la batlla del Baix Aragó // Estudis Universitaris Catalans. 1930. Vol.~15. P.~304--323.}. Уже к этому времени были известны позднейшие нотариально заверенные списки XVI~в. Более полное специальное издание памятника, подготовленное в 1996~г., к сожалению, осталось нам недоступным\footnote{Els Costums d’Orta (1296): Estudi introductori i edició / ed. J. Serrano Daura. Horta de Sant Joan: Ajuntament d’Horta de Sant Joan, 1996. 103 p.}. Поэтому в настоящем исследовании текст «Обычаев Орты» используется по его воспроизведению в сборнике «Исторические законы Каталонии», где опубликован текст памятника и приведены сведения о его рукописной традиции\footnote{Costumbres de Orta. 1296, abril, 16. Orta (Horta de Sant Joan) // Leyes Históricas de Cataluña / dir. J. Serrano Daura. Madrid: Agencia Estatal Boletín Oficial del Estado, 2024. Vol.~2: Apéndices. P.~617--626.}.

С источниковедческой точки зрения «Обычаи Орты» следует рассматривать как результат оформленного в письменном виде компромисса между властью Одена тамплиеров и местной общиной. Инициатива фиксации норм была связана с конфликтом о пределах действия норм «Обычаев Льейды»: жители утверждали, что пользовались и должны пользоваться последними, тогда как представитель Орден признавал за ними лишь отдельные нормы процессуального и коммерческого права.

В преамбуле памятника подчёркнуто, что жители Орты нуждались в \emph{«определённых»} и \emph{«особых» обычаях}\footnote{Costumbres de Orta…: «…certas habere consuetudines ac etiam speciales». P.~617.}. Эта формула показывает, что составители стремились не столько подтвердить прежнюю практику, сколько ограничить возможное произвольное обращение к внешним правовым системам. «Обычаи Орты» зафиксировали переход от разрозненных, частично устных и частично письменных правовых практик использования основных кодификаций обычного права, таких, как «Барселонские обычаи» или «Обычаи Льейды» к утверждённому сеньором локальному письменному своду.

Содержание источника показывает, что власть ордена тамплиеров стремилась совместить контроль над зависимым населением с признанием определённых форм местного самоуправления. Жители Орты сохраняли право пользования природными ресурсами, освобождались от ряда местных поборов, обладали гарантиями личной безопасности и могли иметь присяжных должностных лиц, однако эти должностные лица представлялись командору и утверждались при условии их пригодности. Таким образом, муниципальная автономия была не самостоятельной политической свободой, а санкционированным элементом сеньориального порядка. Судебная власть также оставалась встроенной в юрисдикцию ордена тамплиеров: командор, батлы и иные должностные лица должны были вершить правосудие, но одновременно не могли наказывать без надлежащих свидетелей или иной законной формы доказательства.

Значительное место в «Обычаях Орты» занимали нормы процессуального права. Они регулируют вызов в суд, предъявление жалобы, письменное оформление иска при определённой цене требования, сроки рассмотрения дел, представительство, свидетельские показания, принесение присяги и порядок исполнения судебного решения.

Вероятно, составители «Обычаев Орты» вдохновлялись «Обычаями Льейды», в меньшей степени --- «Барселонскими обычаями», в исключительных случаях обращались к локальному правовому обычаю. Таким образом, «Обычаи Орты» включались в многоуровневую систему источников права Каталонии конца XIII~в.

«Обычаи Орты» представляют собой один из источников для изучения процессов записи местного обычного права, рецепции «Обычаев Льейды», «Барселонских обычаев» на территориях Новой Каталонии конца XIII~в. Их значение состоит не только в фиксации конкретных норм уголовного, гражданского, торгового, семейного и процессуального права, но и в демонстрации самого механизма правового согласования между сеньориальной властью и общиной.

\paragraph*{§2.1.3. «Обычаи Миравета» (1319)}

«Обычаи Миравета» представляют собой один из памятников локального права Новой Каталонии первой трети XIV~в. Миравет был отвоёван у мусульман в 1153~г. Рамоном Беренгером IV и в том же году передан тамплиерам. После упразднения этого ордена часть его владений перешла к госпитальерам, и в 1317~г. Миравет оказался в составе их владений. Составление «Обычаев Миравета» связано с переходом бывших владений тамплиеров под власть госпитальеров ордена святого Иоанна Иерусалимского. Это событие относится к 1319~г.\footnote{Обращение к «Обычаям Миравета», «Обычаям Балагера» и материалам Валь-д’Арана XIV~в. означает сознательный выход за основные хронологические рамки исследования, однако он оправдан его задачами: эти памятники привлекаются не как самостоятельный предмет анализа, а как свидетельства рецепции «Обычаев Льейды». Именно рецепция позволяет проследить дальнейшее функционирование «Обычаев Льейды» как «эталонного» свода обычного права, её адаптацию в иных локальных средах и тем самым точнее определить историческое значение самого памятника.} и связано с деятельностью Мартина Переса де Ороса, кастеляна Ампосты и сеньора Миравета от имени ордена Госпитальеров.

Возникновение текста «Обычаев Миравета» было обусловлено необходимостью устранить неопределённость, существовавшую ещё период сеньории тамплиеров. В Миравете ранее применялись, с одной стороны, «Обычаи Льейды», с другой --- местные нормы неписанного права. Расхождение между этими правовыми массивами вызывало споры между жителями батлии и сеньориальной администрацией относительно того, какое право должно применяться в конкретных случаях. Поэтому «Обычаи Миравета» следует рассматривать как результат кодификации и согласования уже существовавших практик.

Рукописная традиция памятника осложнена тем, что ни каталонский оригинал 1319~г., ни перевод на латинский 1320~г. до настоящего времени не сохранились. Известны лишь более поздние списки: каталонский список XV~в., находящийся в Библиотеке Коломбина в Севилье\footnote{Biblioteca Capitular Colombina (Sevilla). Cod. A A, 141, 19.}, и латинский список 1328~г., хранящийся в Национальном историческом архиве в Мадриде\footnote{AHN. Cód. 944-B.}. Копия на старокаталанском относится к редакции до подтверждения «Обычаев Миравета» в Арле, тогда как её латинский перевод --- к редакции, сложившейся после 1320~г. Это обстоятельство важно потому, что «Обычаи Миравета» дошли до нас на двух языках. Каталонский текст отражает редакцию 1319~г., латинский --- 1320~г. В данном случае речь идёт о существовании памятника в двух языковых версиях уже на раннем этапе его истории.

Более ранняя каталонская версия «Обычаев Миравета» была введена в научный оборот Ф.~Вальсом-и-Табернером, который в 1926~г. опубликовал её сначала в журнале\footnote{\emph{Valls i Taberner F. }Les Costums de la Batllia de Miravet // Revista Jurídica de Catalunya. 1926. Vol.~32. P.~52--76.}, а затем отдельным изданием\footnote{Les Costums de Miravet: editades amb una nota preliminar / ed. F. Valls i Taberner. Barcelona: Estampa de fill de D. Casanovas, 1926.}. При этом он воспроизводил не саму рукопись из Библиотеки Коломбина в Севилье, а её транскрипцию, выполненную в XIX~в. Х.Л.~Вильянуэвой\footnote{\emph{Serrano Daura J. }La concessió dels costums de la batllia de Miravet… P.~295.}. Латинский текст был опубликован Г.~Гало Санчесом в 1915~г. в издании, которое сопровождалось вводными замечаниями о рукописной традиции памятника, соотношении его каталонской и латинской редакций и принципах публикации\footnote{Constituciones Baiulie Mirabeti / ed. G. Sánchez. Madrid: Publicaciones de la Residencia de Estudiantes, 1915.}. Современное критическое издание Дж.~Серрано-Дауры представляет собой специальную публикацию памятника с введением и комментариями\footnote{Els costums de la Batllia de Miravet: estudi introductori i edició / ed. J. Serrano Daura. Gandesa (Tarragona): Consell Comarcal de la Terra Alta, 1999.}. К сожалению, на этапе подготовки настоящей диссертации оно осталось для нас недоступным, поэтому мы пользовались изданием Г.~Гало Санчеса.

Сопоставление каталонской и латинской версий выявляет ряд внешних и редакционных различий. Каталонская рукопись делится на пять книг, а отдельные главы снабжены заголовками, тогда как латинский текст, по наблюдению Дж.~Серрано-Дауры выглядит более ранним по отношению к этой рубрикации\footnote{\emph{Serrano Daura J.} La concessió dels costums de la batllia de Miravet // Revista de dret històric Català. 2015. Vol.~14. P.~295.}. Кроме того, заглавия двух редакций не совпадают: каталонская традиция обозначает памятник как свод обычаев батлии Миравета, тогда как латинская --- как установления. Эти различия не являются чисто формальными: они показывают, что текст памятника существовал в разных документальных режимах --- как локальный свод обычного права и как подтверждённый орденом госпитальеров нормативный акт.

Памятник включает 134 статьи; деление на 5 книг связано исключительно с старокаталанской традицией и, по-видимому, не отражает первоначальной строгой систематизации норм внутри свода. Материал свода охватывает вопросы правового статуса жителей, судебного процесса, долговых отношений, мер и весов, торговли, семейного и наследственного права, уголовных наказаний и полномочий местных должностных лиц.

Особое значение имеет глава о порядке применения источников права. В ней установлено, что при недостаточности собственных норм свода следует обращаться к Конституциям Каталонии, затем к «Барселонским обычаям», а затем к общему праву\footnote{Constituciones Baiulie Mirabeti / ed. G. Sánchez. Madrid: Publicaciones de la Residencia de Estudiantes, 1915. P.~33--34.}. Эта норма позволяет определить место «Обычаев Миравета» в правовой системе Каталонии XIV~в.: они выступали частью многоуровневой системы, при этом их составители за собой более авторитетные дополнительные источники права.

Таким образом, для настоящего исследования «Обычаи Миравета» важны не только как локальный правовой памятник начала XIV~в., но и как свидетельство дальнейшей рецепции и переработки правовой модели, связанной с «Обычаями Льейды». Их рукописная и языковая традиция, как и обстоятельства их составления и подтверждения, показывают, что перед нами текст, возникший в результате последовательной обработки и приспособления уже существующего правового материала к новым условиям. В этом смысле «Обычаи Миравета» позволяют проследить, каким образом правовая традиция «Обычаев Льейды» сохраняла значение за пределами первоначальной среды своего формирования, включалась в иные структуры власти и функционировала в составе более сложного правового пространства. Обращение к этому памятнику, следовательно, даёт возможность судить не только о праве самой батлии Миравета, но и о механизмах передачи, закрепления и адаптации правовых текстов в Каталонии XIII--XIV~вв., что имеет принципиальное значение для проблематики настоящей диссертации.

\subsection*{§3. «Кутюмы Перпиньяна» (1242)}
\addcontentsline{toc}{subsection}{§3. «Кутюмы Перпиньяна» (1242)}

Важнейшим памятником муниципального права Руссильона периода Высокого Средневековья являются «Кутюмы Перпиньяна». Формирование данного свода представляет собой длительный процесс, однако основная кодификация текста относится ко второй половине XII --- началу XIII~в.

Ж.~Массо-Ренье датировал возникновение «Обычаев Перпиньяна» 1190-ми или даже 1170-ми гг., ссылаясь на отсутствие упоминаний о должности городского консула, которая достоверно существовала в XIII~в. и ещё ряд аргументов, позднее опровергнутых Ж.-М.~Карбассом\footnote{\emph{ Carbasse J.-M. }Les anciennes coutumes de Perpignan // Société Agricole Scientifique et Littéraire des Pyrénées-Orientales. 1988. No.~96. P.~181.}. Последний считал наиболее вероятной датировку этого правового компендиума 1242--1246~гг., поскольку, в королевских привилегиях, пожалованных Перпиньяну в 1241~г., «Кутюмы» вовсе не упоминались\footnote{Ibid. P.~181.}.

Авторство документа является коллективным и приписывается деятельности городских консулов и нотаблей, стремившихся зафиксировать привилегии города на фоне укрепления королевской власти Арагонских королей. Одним из доказательств коллективного авторства служит довольно хаотичная внутренняя структура свода, характерная для ранних компиляций обычного права (например, «Барселонских обычаев»)\footnote{Ibid. P.~182.}. В то же время Ж.-М.~Карбасс отметил, что рукопись не содержала никаких следов этой коллективной редактуры (как в случае с «Обычаями Тортосы»), а её лексика коренным образом отлична от компиляций, составленных профессиональными юристами (как, например, в случаях с «Обычаями Льейды», «Кутюмами Арля» и «Кутюмами Нарбонны»)\footnote{Ibid. P.~188.}.

В общей сложности «Кутюмы Перпиньяна» содержат 69 статей, расположенных хаотично и затрагивающих различные области права. Так, всего 34 статьи посвящены гражданскому процессуальному праву, 13 --- торговым обычаям, 6 --- праву уголовному, остальные --- различным вопросам от строительства зданий до положения иудеев, проживавших на территории города. При этом взаимное расположение статей внутри кодекса не позволяет выделить тематические и хронологические блоки.

«Кутюмы Перпиньяна» представляли собой синтез местного обычного права, существовавшего в устной форме, и элементов римского права, проникавшего в регион через «учёное право». Что характерно, уже первая статья «Кутюм Перпиньяна» определяет иерархию источников права на территории города и его округи, формально ограничивая действие собственно каталонского права: \emph{«жителей Перпиньяна следует судить по обычаям города, а там, где нет обычая, --- по римскому праву; а не по „Барселонским обычаям“ или „Вестготской правде“»}\footnote{Ibid. P.~183.}. С другой стороны, по мнению Ж.-М.~Карбасса, 19-я статья «Обычаев Перпиньяна» практически полностью заимствована из 20-й статьи «Барселонских обычаев» в то время, как статьи №~25 и 28 --- из обычаев других городов Южной Франции: Монпелье и Нарбонны, Сен-Жиля, Безье и Каркассона\footnote{Ibid. P.~189.}.

К середине~XIX~в. текст «Кутюмов Перпиньяна» сохранился в составе 3-х картуляриев: «Поместной книги», «Большой Зелёной Книги» и «Малой Зелёной книги»\footnote{\emph{Massot-Reynier J. }Les Coutumes de Perpignan, suivies des usages sur la dîme, des plus anciens priviléges de la ville et de documents complémentaires, publiées en latin et en roman d’après les manuscrits par la société archéologique de Montpellier. Montpellier, 1848. P.~XVII.}. Фундаментальным изданием источника является публикация 1848 года, подготовленная Ж.~Массо-Ренье, в которой «Обычаи Перпиньяна» представлены в виде параллельного текста на латыни и романском диалекте, что делает ее незаменимой для сравнительно-лингвистического анализа. Издание снабжено ценными приложениями, включающими обычаи о десятине и древнейшие привилегии города, что значительно расширяет контекст исследования.

«Кутюмы Перпиньяна» наглядно демонстрируют процесс трансформации устного обычая в писаный закон, характерный для городской среды указанного периода. Памятник позволяет выявить специфические черты каталонского муниципального самоуправления в запиренейских владениях принципата Каталонии и его отличия от аналогичных институтов Южной Франции и Арагона. Кроме того, анализ текста дает возможность проследить проникновение римского права в обычное задолго до окончательной рецепции римского права. Таким образом, «Кутюмы Перпиньяна» частично репрезентативны для компаративного аспекта исследования правовой культуры городского сословия Руссильона и использовались нами для сравнения с каталонскими памятниками.

\subsection*{§4. «Обычаи Тортосы» (1272, 1276--1279)}
\addcontentsline{toc}{subsection}{§4. «Обычаи Тортосы» (1272, 1276--1279)}

«Обычаи Тортосы» были составлены в 1270-е годы, когда в городе произошёл серьёзный конфликт юрисдикций между главой ордена тамплиеров, генуэзцами и представителями графов Барселонских. Менее, чем за 10 лет «Обычаи Тортосы» претерпели две последовательные редакции: в 1272~г. и в 1277--1279~гг. При этом «Обычаи Тортосы» --- один из самых объёмных по числу установлений свод локального обычного права\footnote{\emph{Daileader P. }La coutume dans un pays aux trois religions~: la Catalogne, 1228-1319 // Annales du Midi~: revue archéologique, historique et philologique de la France méridionale. 2006. №~225 (118). P.~372.}. «Обычаи Тортосы» (в обеих редакциях) изначально были разделены на 9 книг, в XIX~в. Б.~Оливером-и-Эстрейером книги были разделены на 142 рубрики и 1350 статей соответственно\footnote{Costums de Tortosa / ed. crít. a cur. de J. Massip i Fonollosa, col. C. Duarte i Montserrat, \emph{M.A.~Massip} i Bonet. Barcelona: Fundació Noguera, 1996. P.~XVIII.}.

Рукописная традиция «Обычаев Тортосы» была представлена 5 сохранившимися рукописями, которые были созданы до XIX~в. Притом, о рукописи А (1272--1273) историки до середины XX~в. не знали, а на момент составления критического издания памятника каталонским историком Дж.~Массипом-и-Фонойоссой одна из этих 5 рукописей была утеряна: 2~рукописи XIII~в. --- A\footnote{ACBEB. Cod. 53.} и B\footnote{ACBEB. Reg. 358. Comú I-113.} (выполнена позднее), 2 рукописи XIV~в. --- C\footnote{Fundació Bartolomé March Servera. Biblioteca. Sign. B96-V2-01.} (c~обширными лакунами) и F (ныне утерянная)\footnote{Ibid. P.~XXXVI.}. В частности, по его мнению, на ныне утерянную рукопись конца XIV--начала XV~вв. опирался Б.~Оливер-и-Эстрейер, известный испанский историк права XIX -- начала XX~в., издатель текста «Обычаев Тортосы». Более поздние рукописи второй редакции памятника --- D и E, согласно стемме, приведённой Дж.~Массипом-и-Фонойоссой в издании 1996~г., были сняты с рукописи C и не являлись самостоятельными\footnote{Ibid. XVIII.}. Первое печатное издание «Обычаев Тортосы», вероятно, было выполнено на основе не дошедших до нас рукописей, составление которых относится к~XV~в.\footnote{Ibid. P.~XVIII.}.

Текст рукописи А, оригинал которой датируется 1272~г. существенно отличается от текста, который сохранился в остальных рукописях. Дж.~Массип-и-Фонойоса убедительно доказал, что в рукопись А, в отличие от последующих, вошёл текст первой редакции «Обычаев Тортосы» 1272~г., а рукописи B, C, D и F -- содержали текст второй версии памятника 1279~г.

История составления и редактирования текста «Обычаев» получила отражение в преамбулах к обеим сохранившимся редакциям интересующего нас памятника. В первой версии (по рукописи А) основной причиной составления «Обычаев Тортосы» назван конфликт городских властей Тортосы и рыцарей-тамплиеров: \emph{«Поскольку спор между братом }\emph{Арнау}\emph{ из }\emph{Кастельноу}\emph{, ордена тамплиеров в Арагоне и Каталонии, братьями из этого ордена, с одной стороны, и горожанами города Тортосы --- с другой, долго рассматривался в римской курии, названный магистр и братья потребовали, чтобы обычаи Тортосы были записаны…»}\footnote{Ibid.: «Con fos pleyt longament estat en la cort de Roma per !frare Amau de Castelnou, Maestre en Aragon s e en Cataluyna, e per los frares de la cavaleria del Temple demanant de la una part, e els ciutadans de la ciutat de Tortosa defenent de l'altra; demanavan lo dit Maestre e els frares que les costumes de Tortosa fossen meses en escrit…». P.~2.}. Преамбула рукописи А в отличие от других версий памятника сообщала о двух нотариях, Петре Тамарите и Петре Эгидии, которым была поручена запись локальных правовых норм.

Составители, перерабатывавшие устные нормы обычаев Тортосы, записывали их на старокаталанском, собственно, на том языке, на котором эти обычаи бытовали. В то же время в отдельных случаях они прибегали к использованию латинского языка.

Однако даже после отбора действовавших норм обычного права, как отмечается в преамбуле, предваряющей текст «Обычаев Тортосы» в рукописях B и С, записанные обычаи оставались \emph{«слишком длинными, неясными и сомнительными во многих местах»}\footnote{Ibid.: «les Costumes romanguessen encara massa longues e escures e duptoses…». P.~3.}, а потому снова были отданы \emph{«господам арбитрам, чтобы те снова могли исправить, изменить, сократить и прояснить, согласно тому, как, по их мнению, было бы правильнее…»}\footnote{Ibid.: «adobar, mudar e abreviar ... a lur discreció». P.~3.}. Таким образом, результат первой попытки записи «Обычаев Тортосы» не удовлетворил обе стороны (и тамплиеров, и городской совет), арбитры установили новую компенсацию в 1000 мараведи (первоначально 2000 мараведи) за неподчинение своим решениям.

Только после повторной правки, редакция «Обычаев Тортосы» была признана состоявшейся. Необходимо отметить, что из второй редакции, созданной 1277--1279~гг., исчезли как указания на имена нотариев-составителей первой редакции, так и описание обстоятельств её составления: эта редакция открывалась рассказом о завоевании Тортосы Рамоном Беренгером IV сразу, за этим повествованием следовало описание работы с текстом «Обычаев Тортосы » в 1275--1277~гг., которое было исключено из текста в ходе правки.

Первая печатная публикация «Обычаев Тортосы» была осуществлена в 1539~г. под редакцией Х.~Амика\footnote{Llibre de les costums generals scrites de la insigne ciutat de Tortosa: ab alguns privilegis, confirmacions e sentencies fahents per a la administració de justicia / ed. J. Amich. Tortosa, 1539. 1 vol.}. После этого в течение нескольких столетий текст не переиздавался. Во второй половине XIX~в. были предприняты две новые публикации: Б.~Оливером-и-Эстельером\footnote{Costums de Tortosa; Libre de les costums generals scrites de la insigne ciutat de Tortosa: texto auténtico de siglo XIII que da a luz nuevamente, ilustrado con observaciones crítico-literarias y un copioso vocabulario / ed. B. Oliver y Esteller. Madrid: Impr. de M. Ginesta, 1881.}, основанное на издании 1539~г. и впервые представившее текст в полном виде, и Р.~Фогетом, выходившее отдельными выпусками и оставшееся незавершённым\footnote{Libre de les costums escrites de la ciutat de Tortosa: texto catalán (auténtico) enriquecido con las variantes de códices anteriores al siglo XV / ed. R. Foguet. Barcelona: C. Verdaguer, 1878--1881.}. В 1912~г. публикация была продолжена Х.~Фогетом Марсалем\footnote{Código de las costumbres escritas de Tortosa a doble texto, traducido al castellano del más auténtico ejemplar catalán / ed. R. Foguet; continuada por J. Foguet Marsal. Tortosa: Impr. Querol, 1912.}. В 1956--1958~гг. Ф.~Х.~Гас Карпио подготовил неполный кастильский перевод, первоначально выходивший как приложение к местному еженедельнику\footnote{Libro de las Costumbres Generales Escritas de la Insigne Ciudad de Tortosa: texto auténtico del siglo XIII / ed. \emph{F.J.~Gas} Carpio. Tortosa: La Voz del Bajo Ebro, 1956--1958.}. В 1972~г. было издано факсимильное воспроизведение рукописи A\footnote{Consuetudines Dertosae: VII centenari 1272--1972 (“Llibre de les costums de Tortosa”) / pról. J. Massip i Fonollosa. Tarragona: Institut d’Estudis Tarraconenses Ramon Berenguer IV, 1972.}, однако, А.~Курто Оммедес оценил его как технически не вполне удачное в виду низкого качества микрофильмов, на основе которых оно было выполнено\footnote{\emph{Curto Homedes A.} Els Costums de Tortosa. Els diferents manuscrits i edicions // Revista de Dret Històric Català. 2023. Vol.~22. P.~109.}. В настоящем исследовании используется наиболее полное и современное издание памятника --- критическая публикация 1996~г., подготовленная Ж. Массипом-и-Фонойосой\footnote{Costums de Tortosa / ed. crít. J. Massip i Fonollosa. Barcelona: Fundació Noguera, 1996.}.

Вопрос о сосуществовании двух языков: классического, универсального, наделенного авторитетом древности, и местного устного языка, который относился к более низкому регистру --- в одном правовом памятнике, представляет особенный интерес. В случае с «Обычаями Тортосы» мы можем наблюдать за тем, как проходил процесс отбора и его редактирования норм обычного права по маргиналиям в рукописи А, в которых сохранились пометки об одобрении или неодобрении епископом Тортосы Арнау и другими духовными лицами, входившими в состав специальной комиссии, отдельных установлений из редакции 1272~г.\footnote{\emph{Calzada González M.A.} El principio ei incumbit probatio qui dicit, non qui negat y su recepción en Les Furs de Valencia y en Les Costums de Tortosa // La prueba y medios de prueba. De Roma al derecho moderno. Madrid: Universidad Rey Juan Carlos, Servicio de Publicaciones, 2000. P.~109--120.}.

Таким образом, «Обычаи Тортосы» занимают в источниковой базе настоящего исследования особое место. Значительный объём памятника, высокая степень сохранности его рукописной традиции и наличие современного критического издания делают его одним из наиболее ценных и информативных источников для изучения правовой культуры и обычного права Каталонии XI--XIII~вв. Именно поэтому «Обычаи Тортосы» позволяют исследовать довольно широкий круг вопросов, связанных с формами правового мышления, механизмами закрепления обычая в письменном тексте и местом городского права в правовом пространстве средневековой Каталонии.

\subsection*{§5. Королевская привилегия «Знатные мужи Барселоны признали…» (Recognoverunt proceres) (1284)}
\addcontentsline{toc}{subsection}{§5. Королевская привилегия «Знатные мужи Барселоны признали…» (Recognoverunt proceres) (1284)}

Королевская привилегия «Знатные мужи Барселоны признали…» (Recognoverunt proceres), дарованная королём Пере~III Великим городу Барселоне 11 января 1284~г., представляет собой фундаментальный памятник каталонского права, знаменующий кульминацию процесса кодификации городских обычаев в XIII~в. Появление этого документа было тесно связано с работой Генеральных кортесов Каталонии и было продиктовано необходимостью укрепления союза короны с городской знатью на фоне внешнеполитических конфликтов с Францией и папством\footnote{Более подробно об историческом контексте см. \emph{Кучумов В.}\emph{А}\emph{.} Кортесы Арагона и Каталонии: генезис и специфика // Проблемы испанской истории [Текст]~: сб / под ред. \emph{Пожарская С.П.}~Москва: Наука, 1992. С.~146--147.}.

Несмотря на то, что формально документ является королевским актом, он был создан в рамках диалога между кортесами и королём. Фактическим автором текста выступил коллектив представителей городской элиты --- \emph{«знатных»} и \emph{«добрых людей»}, действовавших от имени Совета ста. Непосредственную редакцию привилегии осуществила, согласно тексту, группа \emph{«старейшин и знатоков в праве»}\footnote{Recognoverunt proceres. I.: «…antiqui et sapientes in iure…». P.~590.}. Вероятно, в данном случае речь идёт о профессиональных легистах, которые, владея техникой римского права, трансформировали устные запросы горожан и знати в конкретные формулировки на языке права. Роль графа-короля в этом процессе сводилась к функции верховного законодателя, который своим утверждением этой привилегии придал статус писаного закона уже существующим практикам.

Основу источника составили местные неписанные и не включённые в состав «Барселонских обычаев» правовые нормы, регулировавшие торговлю и семейные отношения, а также нормы из более ранних привилегий короля Жауме I Завоевателя (1213--1276).

В дальнейшем эта привилегия 1284~г. была включена в компиляцию, известную под названием «Установления и прочие права Каталонии» (кат.~Constitucions i altres drets de Catalunya). Распространение норм происходило также посредством рецепции права столицы графства Барселонского, когда другие города Каталонии запрашивали привилегию судиться «по обычаю Барселоны», тем самым перенимая положения данного источника.

Рукописная традиция привилегии «Знатные мужи Барселоны признали…» представлена тремя рукописями на латинском и старокаталанском языках. К ней относятся латинский оригинал 1284~г.\footnote{Archivo de la Corona de Aragón. Cancillería. Consell de Cent. Pergaminos. Núm. 5131.}, старокаталанская копия XIV~в. из библиотеки Эскориала\footnote{Real Biblioteca del Monasterio de El Escorial. Sign. O-I-12. Fols. 17--33.} и недатированная копия документа на старокаталанском языке в Историческом архиве Барселоны\footnote{Archivo Histórico de la Ciudad de Barcelona. Pergaminos Municipales. 1A-105.}. Такое соотношение рукописей позволяет учитывать при работе с памятником его раннюю передачу на старокаталанском языке.

Впервые издана в печатном виде привилегия была в 1704~г. в составе официальной компиляции каталонского права\footnote{Constitucions y altres drets de Cathalunya, compilats en virtut del capitol de cort LXXXII de las Corts per la S.C. y R. Majestat del Rey Don Philip IV, Nostre Senyor, celebradas en la Ciutat de Barcelona any M. DCCII. Barcelona: Casa de Joan Pau Martí y Joseph Llopis Estampers, 1704. Vol.~2. Lib. 1. Tit. 13. P.~39--51.}. В XX~в. каталонская версия была издана отдельно Х.~~Ровирой-и-Эрменголем в 1927 г\footnote{Recognoverunt proceres. Versión medieval catalana del privilegio así llamado / ed. J. Rovira. Barcelona: Departamento de Historia del Derecho de la Universidad de Barcelona, 1927.}. Появление этой публикации вызвало особый отклик со стороны видного специалиста по истории каталонского обычного права, Ф.~Вальса-и-Табернера, посвятившего памятнику отдельную статью\footnote{\emph{Valls i Taberner F. }Les consuetuds i franqueses de Barcelona de 1284, o “Recognoverunt proceres” // Revista de Catalunya. 1927. No.~39. P.~248--254.}. Латинский текст позднее был вновь опубликован в составе собрания королевских привилегий Барселоны в 1971~г.\footnote{Privilegios reales concedidos a la ciudad de Barcelona / ed. \emph{A.M.~Arag}ó, \emph{M.M.~Costa}. Barcelona: Archivo de la Corona de Aragón, 1971. Vol.~43. P.~8--17.}. В настоящее время наиболее удобным для работы является воспроизведение обеих версий в сборнике «Исторические законы Каталонии», где латинский оригинал и каталонская версия помещены параллельно и снабжены кратким археографическим комментарием\footnote{Recognoverunt proceres. 1283, enero, 11. Barcelona // Leyes Históricas de Cataluña. Madrid: Agencia Estatal Boletín Oficial del Estado, 2024. Vol.~2: Apéndices. P.~589--617. (В дальнейшем цитируется по этому изданию).}.

Значение этого правового памятника в контексте диссертации определяется не только содержанием закреплённых в нём локальных обычаев и правовых установлений, но и особенностями его текстовой традиции: наличием латинского оригинала и ранней старокаталанской версии. Привилегия «Recognoverunt proceres» может показать, каким образом формировалось каталонское обычное право в момент его письменного закрепления, редакционной обработки и последующего функционирования в городской среде. В рамках темы диссертации он представляет интерес как источник по формам фиксации обычая, по соотношению латинского и народного языков в правовой письменности и по механизмам включения местной правовой традиции в более широкий корпус каталонского права

\subsection*{§6. Актовый материал}
\addcontentsline{toc}{subsection}{§6. Актовый материал}

Отдельную группу источников настоящего исследования составляют грамоты, опубликованные в двух томах серии «Каталонские юридические тексты»: документация о правосудии и разрешении конфликтов в каталонских графствах IX--XI~вв. и аналогичная документация XII в\footnote{Justícia i resolució de conflictes a la Catalunya medieval. Col·lecció diplomàtica. Segles IX--XI / dir. \emph{J.M.~Salrach} i Marès, T. de Montagut i Estragués. Barcelona: Parlament de Catalunya: Departament de Justícia de la Generalitat de Catalunya, 2018. (Textos Jurídics Catalans; 37. Documents; 2); Justícia i resolució de conflictes a la Catalunya medieval. Col·lecció diplomàtica. Segle XII / dir. \emph{J.M.~Salrach} i Marès, T. de Montagut i Estragués. Barcelona, 2024. (Textos Jurídics Catalans; 38. Documents; 2/3).}. Первый из этих сборников включает 557 документов, точнее 552 основных документа и пять дополнительных, относящихся к IX--XI~вв., второй --- 873 документа XII~в. В совокупности оба тома дают корпус из 1430 актов, охватывающих период с 812 по 1199~г. и позволяющих проследить изменения в формах судебной практики, добровольного урегулирования конфликтов на протяжении всей раннесредневековой истории каталонских графств.

Ценность обоих сборников определяется не только широтой документальной базы, но и качеством её научной организации. Составители подробно оговаривают критерии отбора, учитывают текстовую традицию каждого документа, пересматривают датировки и, где это возможно, уточняют место составления. Особое значение имеют развёрнутые регесты, позволяющие видеть не только содержание акта, но и конфликт, процедуру и формы его разрешения. Важную роль играет и справочный аппарат в конце томов: перечни архивов, именные и специальные указатели делают корпус пригодным не только для выборочного чтения, но и для системного анализа правового материала в его нормативном, социальном и институциональном контексте. Существенно также, что в обоих томах в специальных приложениях и комментариях уточняются источники цитат и текстовых заимствований из правовых памятников, к которым сами грамоты прямо или косвенно отсылают. Именно сочетание полноты выборки, тщательной подготовки издания, внимания к правовой терминологии и развитого справочного аппарата делает эти сборники надёжной основой для источниковедческого анализа и автоматизированной обработки текстов грамот, которая была предпринята нами в рамках настоящего исследования.

При работе с этим типом источников в центре внимания оказываются документы, связанные с конкретными спорами, жалобами, судебными заседаниями, переговорами, отказами от притязаний и восстановлением мира. Тем самым они позволяют увидеть право не только как систему предписаний, но и как совокупность действий, жестов, процедур, формул и решений, при помощи которых общество пыталось справиться с конфликтом\footnote{\emph{Варьяш}\emph{ }\emph{О}\emph{.}\emph{И}\emph{.} Проблемы социализации индивида // \emph{Она}\emph{ }\emph{же}\emph{.} Пиренейские тетради: право, общество, власть и человек в средние века. Москва: Наука, 2006. С.~33.}. Именно поэтому такие грамоты особенно важны для реконструкции правового пространства: они фиксируют не отвлечённую норму, а момент её применения, обсуждения, обхода или переопределения.

Типология документов, вошедших в оба сборника, чрезвычайно разнообразна. Среди них встречаются акты судебных заседаний, судебные решения и приговоры, соглашения о мире и урегулировании, записи свидетельских показаний, прошения, экспертные заключения о порядке ведения процесса, обязательства явиться в суд и принять решение, документы о предоставлении гарантий и поручителей, акты осмотра и разграничения спорных территорий, документы о взыскании штрафов, а также многочисленные грамоты, фиксирующие отказ от прав, признание правоты другой стороны или прекращение спора. Уже одно это разнообразие показывает, что речь идёт не о едином типе источников, а о сложном комплексе актового материала, внутри которого каждый привлекаемый документ требует самостоятельной оценки.

Работа с таким материалом, как утверждает И.И.~Варьяш, неизбежно строится на балансе «между казусом и типологией»\footnote{\emph{Варьяш}\emph{ И.И.} Сарацины под властью арагонских королей: исследование правового пространства XIV~в. Санкт-Петербург: Евразия Clio, 2016. С.~11.}. Взятая отдельно грамота часто кажется единичным и даже случайным свидетельством: спор о земле, отказ от притязаний, жалоба на насилие, признание чужого права, соглашение между родственниками или соседями. Однако в совокупности такие казусы обнаруживают устойчивые формы правовой реакции. Они показывают, какие конфликты становились предметом письменной фиксации, кто мог инициировать их разрешение, какие лица выступали посредниками, арбитрами, судьями или свидетелями, какие формулы использовались для прекращения спора и каким образом письменный акт превращал социальное напряжение в юридически признанный результат.

В грамотах постоянно присутствует память о писаном законе, прежде всего о вестготской правовой традиции, но вместе с тем сохраняется значение переговоров, судебных решений, поручительства, отказа от притязаний и восстановления мира. Для XI--XII~вв. особенно заметно, что письменная фиксация конфликта не всегда служила только подтверждением судебной победы одной стороны.

При использовании этих документов необходимо учитывать и способ их отбора. Составители сознательно исключили из корпуса законодательные тексты и общие предписания, а также документы, формально связанные с судебной процедурой, но не отражающие конфликта, например акты подтверждения завещаний\footnote{Justícia i resolució de conflictes a la Catalunya medieval. Col·lecció diplomàtica. Segle XII / dir. \emph{J.M.~Salrach} i Marès, T. de Montagut i Estragués. Barcelona, 2024. P.~XIV.}. Такой выбор важен для источниковедческой оценки: перед исследователем находится не вся правовая письменность каталонских графств, а именно та её часть, где право проявляется через спор и его разрешение. Следовательно, корпус особенно информативен для изучения конфликтности, судебной процедуры, форм примирения и языка правовых притязаний, но менее пригоден для прямой реконструкции всего объёма действовавших норм.

Особую осторожность вызывают документы отказа от прав, признания и «умиротворения». В IX--XI~вв. подобные акты постепенно превращались в устойчивый формуляр, и далеко не всегда за словами об отказе или признании можно с уверенностью увидеть реальный конфликт. Иногда текст сохраняет лишь правовой результат, тогда как ход спора, его причины и соотношение сил между сторонами остаются скрытыми. Поэтому такие грамоты нельзя автоматически читать как судебные дела. Их следует анализировать с учётом формулы, контекста, участников, предмета спора и, по возможности, связанных документов. Только тогда можно определить, идёт ли речь о завершении конфликта, о добровольной передаче права или о стереотипной записи.

Не менее важно, что данные грамоты сохраняют «голос» конфликта не непосредственно, а через письмо. Даже когда документ излагает жалобу, признание или показание свидетеля, перед нами не устная речь участника, а её письменная обработка. Писец, судья, нотариальный или монастырский круг, в котором создавался акт, задавали форму изложения, лексику и правовую рамку. Поэтому грамота сообщает не только о самом конфликте, но и о том, как данный конфликт был понят, классифицирован и приведён в приемлемую с точки зрения права форму. Иначе говоря, источник фиксирует не только событие, но и способ его правового осмысления.

Внешняя критика этих грамот предполагает также внимание к месту их хранения и публикации. Большинство документов дошло в составе картуляриев, монастырских и церковных архивов, собраний документов и позднейших копий. Это означает, что состав используемого комплекса источников неоднороден: лучше представлены те споры, которые по тем или иным причинам оказались важны для учреждений, заинтересованных в письменном закреплении своих прав. Поэтому документы особенно хорошо освещают имущественные конфликты, споры о земле, юрисдикции, зависимых людях, церковных владениях, границах и обязательствах. Напротив, значительная часть повседневной правовой практики, не требовавшей письменного оформления или не сохранившейся в архивах крупных институций, остаётся за пределами прямого наблюдения.

Ключевая особенность этого комплекса источников состоит в том, что он позволяет проследить изменение языка права. В документах IX--XI~вв. особенно заметна связь судебной практики с вестготской традицией, письменными доказательствами и авторитетом судьи. В XI~в. возрастает значение договорных форм разрешения конфликтов. В XII~в. документы становятся более пространными, насыщенными специализированной лексикой и правовой аргументацией и даже прямой речью. Составители второго тома прямо подчёркивают, что грамоты XII~в. длиннее и содержательно насыщеннее, а их язык и аргументация требуют от историка более внимательного чтения\footnote{Justícia i resolució de conflictes a la Catalunya medieval. Col·lecció diplomàtica. Segle XII / dir. \emph{J.M.~Salrach} i Marès, T. de Montagut i Estragués. Barcelona, 2024. P.~XIII.}.

Для темы настоящей диссертации этот материал имеет дополнительное значение. Записи городских обычаев XII--XIV~вв. наследовали более ранним формам фиксации споров, местных норм, писанного и неписанного права: жалобам, соглашениям, отказам от притязаний, свидетельским показаниям, судебным решениям, актам разграничения и гарантиям явки в суд. Поэтому документы из этих двух сборников важны не только как самостоятельный источник по истории правосудия. Именно они могут показать, из какой практической среды постепенно вырастала потребность в письменной фиксации обычаев, процедур и устойчивых форм правового поведения.

С одной стороны, конфликты, о которых идёт речь в грамотах, уже были зафиксированы письменно, а значит так или иначе находились в поле зрения местных светских или церковных властей. Эти сборники не следует рассматривать как полное собрание всех документов, связанных с правом и судом. Сами составители подчёркивают, что абсолютная полнота здесь едва ли достижима: документы такого рода не сосредоточены в одном архиве или в одном типе рукописных книг, а рассеяны по дипломатариям, картуляриям, архивным собраниям, приложениям к монографиям и статьям\footnote{Justícia i resolució de conflictes a la Catalunya medieval. Col·lecció diplomàtica. Segle XII / dir. \emph{J.M.~Salrach} i Marès, T. de Montagut i Estragués. Barcelona, 2024. P.~XIII.}.Однако именно в этом состоит их источниковедческая ценность. Исключительно актовый материал позволяет взглянуть на правовую культуру сквозь призму конкретного спора: увидеть жалобу, спор, процедуру, клятву, ордалию, свидетельство, соглашение, отказ, наказание и восстановление нарушенного порядка\footnote{\emph{Варьяш}\emph{ И.И.} Сарацины под властью арагонских королей: исследование правового пространства XIV~в. Санкт-Петербург: Евразия Clio, 2016. С.~13.}. В совокупности эти документы дают возможность изучать не только отдельные правовые казусы, но и саму правовую культуру каталонских графств IX--XII~вв., в которой обычай, власть, насилие, конкретный человек и письменная фиксация права были связаны между собой значительно теснее, чем это позволял бы установить один лишь нормативный материал.

Поэтому комплекс источников носит создавался а посредством выявления и отбора грамот из уже известных старых и новых публикаций. Это выборка, сформированная на основе содержательных критериев, по которым проводился отбор документов для публикации. При этом её объём и характер позволяют считать её достаточно репрезентативной для изучения правовой культуры каталонских графств IX--XII~вв.

Такое решение продиктовано не только исследовательской задачей, но и практическими пределами работы: чрезвычайное богатство каталонского документального наследия делает невозможным его исчерпывающий охват в рамках одного исследования. Как показал П.~Боннасси, именно сопоставление картуляриев с собраниями подлинников в Каталонии позволяет по-новому поставить вопрос о репрезентативности документальных комплексов\footnote{\emph{Bonnassie P. }La Catalogne du milieu du Xe à la fin du XIe siècle. Croissance et mutations d’une société. Toulouse, 1975. T.~1. P.~23--25.}.

В свою очередь, И.С.~Филиппов справедливо настаивал на том, что оценка любой группы источников требует соотнесения её с историей архива и с типологически сходным, но лучше сохранившимся материалом\footnote{\emph{Филиппов И.С.} Рукописи не горят: проблемы реконструкции погибших архивов и возможности расширения базы источников по истории Средневековья // Как мы пишем историю? М., 2013. С.~199--202}. Поэтому опора на уже опубликованные и тематически отобранные документы не отменяет исследовательской ценности корпуса, но задаёт границы выводов: они неизбежно зависят от степени сохранности, принципов отбора и истории публикаций источников.

Таким образом, источниковая база настоящей диссертации отличается одновременно разнородностью и внутренней связностью. В неё входят нормативные памятники различного происхождения, объёма и степени сохранности --- от «Барселонских обычаев» и связанных с ними сводов городского права до локальных кодификаций XIII--XIV~вв., а также актовый материал, позволяющий видеть право в действии. Эта совокупность источников показывает, что правовая культура Каталонии XI--XIII~вв. складывалась не в пределах одного текста или одной правовой традиции, а на пересечении нескольких уровней: обычая, его письменной фиксации, перевода, редакционной переработки, рецепции в новых локальных средах и практического применения в судебной и внесудебной сфере. Именно поэтому для настоящего исследования принципиально важно учитывать не только содержание отдельных норм, но и язык памятников, историю их рукописной передачи, характер публикаций, место каждого текста в системе источников права и формы его последующего функционирования.

Для темы диссертации значение имеют не только наиболее известные и авторитетные своды, но и их переводы, позднейшие редакции, локальные переработки и актовый контекст, в котором право приобретало действенность. Нормативный текст и документ в данном случае не противопоставлены друг другу: напротив, только в их соотнесении становится возможным увидеть, каким образом в средневековой Каталонии право осмыслялось, записывалось, переводилось, приспосабливалось к новым условиям и включалось в повседневные практики. Поэтому источниковедческий анализ в рамках настоящей работы направлен не только на установление происхождения и статуса отдельных памятников, но и на выявление тех механизмов правовой культуры, которые связывали между собой язык, норму, институт, социальное действие и письменную форму.

Особое место в этой системе занимают вспомогательные источники, привлекаемые для сравнительного анализа: тексты Священного Писания, \emph{Corpus }\emph{Iuris}\emph{ Civilis}, V книга «Этимологий» Исидора Севильского, «Вестготская правда» и «Составленные Петром извлечения из римских законов» как интеллектуальный и правовой контекст. Обращение к ним позволяет точнее определить характер заимствований, различить прямое цитирование, устойчивые формулы, понятийные переклички и более глубокие следы правовой и книжной традиции. Тем самым вспомогательные источники дают возможность увидеть исследуемые тексты не изолированно, а в более широком пространстве правовой культуры, где местный обычай, учёное право, библейская и энциклопедическая книжность, а также позднеантичное законодательное наследие вступали в сложное и продуктивное взаимодействие.
